%\documentclass[manuscript,blind]{geophysics}
%\documentclass[paper,twocolumn,twoside]{geophysics}
\documentclass[manuscript,noblind]{geophysics}

% An example of defining macros
\newcommand{\rs}[1]{\mathstrut\mbox{\scriptsize\rm #1}}
\newcommand{\rr}[1]{\mbox{\rm #1}}

%% Extra packages
\usepackage{amsmath}
\usepackage{amssymb}
\usepackage{amsthm}
\usepackage{nicefrac}
\usepackage{array}
\usepackage[linesnumbered,lined,ruled,commentsnumbered]{algorithm2e}

% replace the term "Input" by "Initialization" in algorithms
\SetKwInOut{Input}{Initialization}

%% For creating theorems, corollaries, lemmas and definitions
%% source: https://www.overleaf.com/learn/latex/Theorems_and_proofs
%\newtheorem{theorem}{Theorem}[section]
\newtheorem{theorem}{Theorem}
%\newtheorem{corollary}{Corollary}[theorem]
%\newtheorem{proposition}{Proposition}[theorem]
\newtheorem{proposition}{Proposition}
%\newtheorem{lemma}{Lemma}[theorem]
%\newtheorem{lemma}{Lemma}
%\theoremstyle{definition}%
\newtheorem{definition}{Definition}

%% Define the QED symbol as a black square
\renewcommand\qedsymbol{$\blacksquare$}

% for creating tables
%\usepackage{booktabs}
%\usepackage{array}

\begin{document}
\title{Planar equivalent layers for potential fields: Part 1 -- Theoretical insights}

\renewcommand{\thefootnote}{\fnsymbol{footnote}} 

\ms{GEO-2024XXXX} % manuscript number

\address{
	\footnotemark[1] Corresponding author: vanderlei@on.br
	\footnotemark[2] Observat\'orio Nacional, Rio de Janeiro, Brazil\\
	\footnotemark[3] Universidade do Estado do Rio de Janeiro, Rio de Janeiro, Brazil\\
	\footnotemark[4] TU Bergakademie Freiberg, Institute of Geophysics and Geoinformatics, Freiberg, Germany
}
\author{Vanderlei C. Oliveira Jr.\footnotemark[2]\footnotemark[1], Christian Gerhards\footnotemark[4]
Val\'eria C. F. Barbosa\footnotemark[2], Andr\'e L. A. Reis \footnotemark[2]\footnotemark[3],
Edson F. Luza\footnotemark[2], Xinpeng Huang\footnotemark[4]}

%\footer{Example}
%\lefthead{Oliveira Jr. et al.}
\righthead{Planar equivalent layers: Part 1}

\maketitle

% Main body
\begin{abstract}
	A comprehensive theoretical framework is developed for potential-field transformations 
	based on planar equivalent layers, a technique empirically used in applied geophysics since 
	the 1960s. The formulation replaces the widely accepted equivalent-layer principle -- primarily 
	founded on empirical generalizations -- with a rigorous theoretical foundation. The planar 
	equivalent-layer theorem establishes the existence of continuous planar layers capable of 
	reproducing potential-field data within a predefined error level, even when the layer 
	intersects the true sources. Closed-form expressions for planar equivalent layers are derived 
	by modifying the classical solutions of Dirichlet’s and Neumann’s problems, revealing that 
	these layers possess tangible physical meaning related to boundary-value problems in potential 
	theory. The theoretical framework also demonstrates the mathematical equivalence between 
	potential-field transformations computed via the equivalent-layer technique and those obtained 
	through Fourier-domain approaches. Distinct classes of equivalent layers are formally defined, 
	together with their relationships to the observed data, true sources, and potential-field 
	transformations. The resulting formulation provides a rigorous explanation for the 
	long-standing success of the equivalent-layer technique in potential-field analysis.
\end{abstract}
\section{Introduction}

The equivalent-layer technique was introduced during the 1960s and has become a versatile
tool for processing and interpreting potential-field data in applied geophysics. 
It consists in approximating a discrete set of potential-field data 
by a linear combination of harmonic functions representing the potential field 
due to a discrete set of virtual sources (e.g., monopoles, dipoles, prisms).
These virtual sources are called ``equivalent sources'' and are commonly disposed on a 
surface forming a discrete ``equivalent layer'', which is ``equivalent" 
in the sense that it satisfactorily reproduces the observed potential-field data.
The general approach involves estimating the constant physical property of the equivalent 
sources (e.g., density, mass, isotropic magnetic susceptibility, magnetization intensity, 
magnetic moment intensity) by numerically solving a constrained linear inverse problem.
The technique can be used to accomplish various data processing steps, such as interpolation, 
upward (or downward) continuation, gradient data processing, merging multiple data sets, 
estimating density models from gravity data, 
computing the amplitude of the anomalous magnetic field, 
transforming total-field anomalies into field components, 
as well as estimating the magnetization direction and computing the reduction to the pole
of magnetic data. 
Further details on practical and computational aspects of the equivalent-layer technique 
can be found in the recent review by \cite{oliveirajr_etal2023}.

Several equivalent-layer methods follow the pioneer approach presented
by \cite{dampney1969} and rely on the numerical solution of the 
``upward continuation integral'' \citep[e.g.,][]{skeels1947, peters1949, henderson-zietz1949,
dean1958, roy1962, bhattacharyya1967, henderson1970}, which is the solution of the
``Dirichlet’s problem'' or the ``first boundary value problem of potential
theory'' on a plane located above the sources \cite[][ p. 236]{kellogg1967}. 
From the physical point of view, this numerical solution gives rise to a predicted potential field 
analogous to the vertical component of the gravitational attraction due to a layer of monopoles
or the magnetic scalar potential due to a layer of vertically-magnetized dipoles.
This approach provides the theoretical background for computing the interpolation or 
upward/downward continuation of any potential-field data.
However, other field transformations that have been successfully used in practice, 
such as joint processing of gravity-gradient data 
\citep[e.g.,][]{barnes_lumley2011} or conversion of gravity disturbance data into the 
gravity-gradient tensor \citep[e.g.,][]{piauilino_etal2019_seg, oliveirajr_etal2023}, are 
empirical generalizations often assumed to be valid despite the lack of mathematical proof.

Most of the well-known equivalent-layer methods for magnetic data processing are 
empirical generalizations of the ``Neumann's problem'' or the 
``second boundary value problem of potential theory'' \cite[][ p. 246]{kellogg1967}.
These methods usually presume, without any proof, the existence of an equivalent layer of 
dipoles located above the sources that can be used to reduce magnetic data to the pole 
\citep[e.g.,][]{emilia1973, bhattacharyya_chan1977, silva1986, leao-silva_1989, oliveirajr_etal2013, 
li_etal2014, oliveirajr_etal2015, ribeiro_filho_etal2020, reis_etal2020, li-etal2021, takahashi-etal2022, piauilino_etal2025}. 
Up to now, far too little attention has been paid to the theoretical aspects related to the 
existence of a planar equivalent layer of dipoles or its relationship with a boundary value
problem of potential theory in applied geophysics.
By following different approaches in the wavenumber domain, 
\cite{pedersen1991} and \cite{li_etal2014} proved the existence of a
planar equivalent layer of dipoles with all-positive
magnetic moment distribution that exactly reproduces the approximated total-field anomaly
\citep[e.g.,][p. 179]{blakely1996} caused by sources having a purely induced magnetization
in the vertical direction.
Their approach, however, is valid only for the case in which 
the total-field anomaly data rely on a plane.
\cite{reis_etal2020} generalized the works of \cite{pedersen1991} and \cite{li_etal2014}
to prove that this planar equivalent layer exists for the case in which the dipoles have 
the same arbitrary uniform magnetization direction of the sources, regardless of whether 
the sources magnetization is purely induced or not.
All those studies assume that the equivalent layer is defined in a source-free region, where  
Laplace’s equation is obeyed. Consequently, all those studies presume that the equivalent 
layer does not cross the true sources.

The seminal work of \cite{gunn1975} was the pioneer in showing a complete description of 
potential-field transformations in the Fourier domain. That author has also formulated such 
transformations in terms of planar equivalent layers of monopoles and dipoles.
However, while presuming the existence of planar layers, \cite{gunn1975} did not establish 
any relationship with boundary-value problems in potential theory and did not take into 
account situations in which they intersect the true sources.

There have been no attempts to establish the theoretical background explaining, for example,
the existence of a planar layer of monopoles that simultaneously fit the gravitational  
potential and its spatial derivatives of first and second orders; 
the existence of a planar layer of dipoles that simultaneously fit the magnetic scalar 
potential and induction field components due to arbitrarily magnetized sources; or the 
existence of equivalent layers that cross the sources.
To the authors knowledge, all equivalent-layer methods found in the literature 
consider that the layer is located above the sources; 
a condition that cannot be ensured in practical situations with field data produced by 
anomaly sources having undisclosed depths.
The absence of precise depth information about the true sources, however, 
does not prevent the real-world application of the equivalent-layer technique.
Existing literature contains numerous examples of its use with a wide range of real-data 
applications, from exploration to satellite geophysics \citep[e.g.,][]{gonzalez_etal2022}.

Here, we present a systematic theoretical 
discussion on the existence of planar equivalent layers 
that reproduce the gravitational and magnetic induction fields due to 3D sources, 
even for the case in which the layer crosses them.
We limit ourselves to the case in which the magnetic
sources have a total magnetization with constant direction, but spatially-variable intensity.
We formulate the existence of continuous planar layers as the 
``planar equivalent-layer theorem'', which we mathematically prove by modifying the 
solutions of the Dirichlet's and Neumann's problems 
on planar surfaces that may cross the sources. 
Our results fill the theoretical gaps in most of the existing equivalent-layer methods, 
including those presenting alternative formulations 
\citep[e.g.,][]{cordell1992, guspi_novara2009}.
We also establish how a particular equivalent layer can be used to perform a particular 
potential-field transformation.
Besides, we define the theoretical link between potential-field transformations made via 
equivalent-layer technique and those directly computed from the Fourier transform of the 
observed potential-field data. 

We address theoretical aspects on the existence of planar equivalent layers and their 
relationship with potential-field transformations with the purpose of replacing with the 
planar equivalent-layer theorem the widely accepted ``equivalent-layer principle''. 
This principle refers to the 
empirical generalizations by which most of the studies found in the geophysical literature 
justify the existence of equivalent layers allowing computing potential-field transformations. 
Rather than a totally new application for the equivalent-layer technique, the main goal here 
is to explain or shed new light on why the technique has long been successfully applied to 
perform potential-field transformations. 
The present work opens the door for future investigations about, for example,  
(i) the analysis of discretization and truncation effects applied to continuous equivalent 
layers; 
(ii) the generalization of the reduction to the pole (RTP) 
to the case in which the sources have a nonuniform total-magnetization direction;
(iii) impact on potential-field transformations of planar non-null physical-property 
distributions that produce no field, that is, the annihilators \cite[][ p. 66]{parker1994} or 
silent sources \citep{baratchart2013, baratchart_etal2021}; and also 
(iv) the relationship between particular planar equivalent layers and the true 3D sources.

While the focus of this work is on the theoretical foundations of the equivalent-layer technique, 
we do not intend to discuss the geophysical implications or the practical computational implementation 
of the theoretical framework presented here. 
To this purpose, we refer the readers to the companion paper \citep{oliveirajr_etal2026_2}. 
In the companion paper, we discuss how the present theoretical developments shown here can be 
used to explain some important theoretical gaps in the equivalent-layer technique. 
We also numerically illustrate how different equivalent layers can be used 
to perform the same potential field transformations and also show an application to field data. 
\section{General background}
\label{sec:background}

\subsection{Preliminary notational and geometric definitions}

The main geometric objects used along the present work are shown in Figure \ref{fig:geometry}.
We consider a topocentric Cartesian coordinate system with axes 
$x$, $y$, and $z$ pointing to north, east and down, respectively. 
We also assume to have a set of 3D gravity or magnetic sources with top at the vertical coordinate $z^{t}_{0}$, 
located in the subsurface over $\overline{\mathcal{V}} = \mathcal{V} \, \cup \, \partial \mathcal{V}$, where 
$\mathcal{V} \subset \mathbb{R}^{3}$ is a \textit{regular region} with boundary
$\partial \mathcal{V}$. 
Some important properties of $\mathcal{V}$ in this case is that it is an open connected set of points 
and $\partial \mathcal{V}$ is a closed and compact surface free of double points \citep[][ p. 22]{freeden_gerhards2012}.
Generic points inside the sources are represented by $\mathbf{r}' = (x', y', z')^{\top} \in \overline{\mathcal{V}}$.

We also consider to have potential field observations at generic points 
$\mathbf{r} = (x, y, z)^{\top} \notin \overline{\mathcal{V}}$ or at a finite set of points 
$\mathbf{r}_{i} = (x_{i}, y_{i}, z_{i})^{\top} \notin \overline{\mathcal{V}}$, $i \in \{ 1, \dots, N \}$, 
located above the sources, i.e., with vertical coordinates satisfying $\forall z < z^{t}_{0}$ and 
$\forall z_{i} < z^{t}_{0}$. These points are conveniently called \textit{observation points}.

Specific horizontal surfaces are defined by $\mathcal{S}$, with a particular superscript symbol. 
The horizontal surface $\mathcal{S}''$, for example, is defined at a vertical coordinate $z''_{0}$ below the 
observation points, so that $\forall z < z''_{0}$ and $\forall z_{i} < z''_{0}$. 
The subscript ``$0$'' indicates that the corresponding vertical coordinate is constant. 
Points and elements of area on a given horizontal surface are defined by using the corresponding superscript symbol. 
For example, we define the element of area $d\mathcal{S}'' = dx''dy''$, generic points $\mathbf{r}'' = (x'', y'', z''_{0})^{\top} \in \mathcal{S}''$ 
and also a finite set of points $\mathbf{r}''_{j} = (x''_{j}, y''_{j}, z''_{0})^{\top} \in \mathcal{S}''$,
$j \in \{ 1, \dots, M \}$.

Finally, all Fourier transforms computed here are defined on horizontal surfaces, with respect to the 
spatial coordinates along the horizontal directions. For example, $\mathcal{F} [ U(\mathbf{r}^{\dagger}) ]$
indicates the two-dimensional Fourier transform of a scalar function $U(\mathbf{r}^{\dagger})$ with respect to
the spatial coordinates $x$ and $y$, on a horizontal surface $\mathcal{S}^{\dagger}$ formed by points  
$\mathbf{r}^{\dagger} = (x^{\dagger}, y^{\dagger}, z^{\dagger}_{0})^{\top}$, i.e. by $\mathbf{r}^{\dagger} \in \mathcal{S}^{\dagger}$.
Besides, the dependence of $\mathcal{F} [ U(\mathbf{r}^{\dagger}) ]$ on the wavenumbers $k_{x}$ and $k_{y}$ associated with 
$x$ and $y$ directions, respectively, is conveniently omitted to simplify the notation.

\subsection{Potential fields due to 3D sources}
\label{sec:potential_fields_3D_sources}

We define the potential field $d^{o}(\mathbf{r})$ produced by the 3D sources at observation points 
$\mathbf{r}$ (Figure \ref{fig:geometry}) by
\begin{equation}
	d^{o}(\mathbf{r}) = D_{d} \, U(\mathbf{r}) \: , 
	\quad \mathbf{r} \notin \overline{\mathcal{V}} \: ,
	\label{eq:observed-data-space}
\end{equation}
where $D_{d}$ is a space-domain differential operator acting on the potential
\begin{equation}
	U(\mathbf{r}) = \kappa \iiint\limits_{\mathcal{V}}
	\sigma(\mathbf{r}') \: \frac{1}{\| \mathbf{r} - \mathbf{r}'\|}
	\: d\mathcal{V}' \: , 
	\quad \mathbf{r} \notin \overline{\mathcal{V}} \: ,
	\label{eq:potential-U-space}
\end{equation}
$d\mathcal{V}' = dx'dy'dz'$ represents a volume element;
$\kappa$ is a proportionality constant;
$\tfrac{1}{\| \mathbf{r} - \mathbf{r}' \|}$ is the inverse distance function
between the observation point $\mathbf{r}$ and a point $\mathbf{r}' \in \mathcal{V}$
and $\sigma(\mathbf{r}')$ is a scalar function defined at points 
$\mathbf{r}' = (x', y', z')^{\top} \in \overline{\mathcal{V}}$ 
as the density (contrast) or total-magnetization intensity within the sources.
We implicitly assume that $\sigma(\mathbf{r}) = 0$ outside $\overline{\mathcal{V}}$.
To insure that the partial derivatives of first and second orders of $U$ 
exist at all points in the $\mathbb{R}^{3}$ space, including those within the sources,
we impose that $\sigma(\mathbf{r}')$ is a piecewise continuous, integrable and bounded function
that satisfies some Hölder condition \citep[][ p. 152]{kellogg1967} at all points
$\mathbf{r}' \in \overline{\mathcal{V}}$.

The Fourier transform of the potential field $d^{o}$ (equation \ref{eq:observed-data-space}) 
on a horizontal surface above the sources is given by \citep[e.g.,][]{bhattacharyya1967,gunn1975}:
\begin{equation}
	\mathcal{F} [ d^{o}(\mathbf{r}^{\dagger}) ] =
	\alpha_{d} \, \mathcal{F} [ U(\mathbf{r}^{\dagger}) ] \: ,
	\quad \mathbf{r}^{\dagger} \in \mathcal{S}^{\dagger} \: ,
	\quad z^{\dagger}_{0} < z^{t}_{0} \: ,
	\label{eq:observed-data-Fourier}
\end{equation}
where 
\begin{equation}
	\mathcal{F} [ U(\mathbf{r}^{\dagger}) ] = 2 \pi \kappa 
	\int\limits_{z^{t}_{0}}^{\infty}
	\mathcal{F} [ \sigma(\mathbf{r}') ] \:
	\frac{e^{- k \, (z' - z^{\dagger}_{0})}}{k}
	\: dz' \: , 
	\quad \mathbf{r}^{\dagger} \in \mathcal{S}^{\dagger} \: ,
	\quad z^{\dagger}_{0} < z^{t}_{0} \: ,
	\label{eq:potential-U-Fourier}
\end{equation}
is the Fourier transform of the Newtonian potential 
$U$ (equation \ref{eq:potential-U-space}) on $\mathcal{S}^{\dagger}$, with $k = \sqrt{k_{x}^{2} + k_{y}^{2}}$, and
\begin{equation}
	\mathcal{F} [ \sigma(\mathbf{r}') ] = 
	\int\limits_{-\infty}^{\infty}\int\limits_{-\infty}^{\infty}
	\sigma(\mathbf{r}') \: e^{-\mathtt{i} \, (k_{x} \, x' + k_{y} \, y')}
	\: dx'dy'
	\label{eq:sigma-Fourier}
\end{equation}
is the Fourier transform of $\sigma(\mathbf{r}')$ 
(equation \ref{eq:potential-U-space}) on the horizontal plane at the vertical
coordinate $z'$ and $\mathtt{i}$ represents the imaginary unit.
In equation \ref{eq:observed-data-Fourier}, $\alpha_{d}$ is a function of the 
wavenumbers $k_{x}$ and $k_{y}$ and represents the differential operator 
$D_{d}$ (equation \ref{eq:observed-data-space}) in the Fourier domain.

The differential operators $D_{d}$ and $\alpha_{d}$ associated with the potential fields commonly found 
in geophysical applications (equations \ref{eq:observed-data-space} and \ref{eq:observed-data-Fourier}) 
can be defined as follows:
\begin{equation}
	D(\ell, \hat{\mathbf{u}}) = 
	\left( \hat{\mathbf{u}}^{\top} \nabla \right)^{\ell} \Leftrightarrow
	\alpha(\ell, \hat{\mathbf{u}}) = \left( \hat{\mathbf{u}}^{\top} \mathbf{k} \right)^{\ell} \: ,
	\quad \ell \in \{ 0, 1, 2 \} \: ,
	\label{eq:diff-operators-generic}
\end{equation}
where $\hat{\mathbf{u}}$ is a constant unit vector,
$\nabla$ is the vector differential operator nabla \citep[e.g.,][ p. 179]{kaplan2002} and 
$\mathbf{k} = ( \mathtt{i} \, k_{x} \, , \, \mathtt{i} \, k_{y} \, , \, k)^{\top}$
\citep{pedersen1978}.
The differential operator $D$ (equation \ref{eq:diff-operators-generic})
has a symbolic form based on that used by \citet[][ p. 179]{kaplan2002}, for example, where the suggested
dot product and exponentiation actually leads to a differentiation.

By combining successive applications of $D(\ell, \hat{\mathbf{u}})$ (equation \ref{eq:diff-operators-generic})
with properly chosen $\ell$ and $\hat{\mathbf{u}}$, we may define differential operators that 
result in all potential fields $d^{o}$ commonly used in Geophysics.
Note that $D(0, \hat{\mathbf{u}}) = \alpha(0, \hat{\mathbf{u}}) = 1$ produces no differentiation.
It can be verified, for example, that using $\ell = 1$ and a constant unit vector $\hat{\mathbf{u}}$
in equation \ref{eq:diff-operators-generic} leads to the following directional derivative 
of first order
\begin{equation}
	D(1, \hat{\mathbf{u}}) \, U(\mathbf{r}) = \nabla U(\mathbf{r})^{\top} \hat{\mathbf{u}}
	\label{eq:potential-U-oblique-1st-deriv-space}
\end{equation}
where
\begin{equation}
	\nabla U(\mathbf{r}) = \left( 
	\tfrac{\partial}{\partial x} U(\mathbf{r}),
	\tfrac{\partial}{\partial y} U(\mathbf{r}),
	\tfrac{\partial}{\partial z} U(\mathbf{r}) \right)^{\top}
	\label{eq:potential-U-grad-space}
\end{equation}
is the gradient of $U(\mathbf{r})$ (equation \ref{eq:potential-U-space}), with elements
given by the partial derivatives of first order
\begin{equation}
	\tfrac{\partial}{\partial u} U(\mathbf{r}) = \kappa \iiint\limits_{\mathcal{V}}
	\sigma(\mathbf{r}') \: 
	\tfrac{\partial}{\partial u} \left( \frac{1}{\| \mathbf{r} - \mathbf{r}'\|} \right)
	\: d\mathcal{V}' \: , 
	\quad \mathbf{r} \notin \overline{\mathcal{V}} \: , \quad u \in \{ x, y, z \} \: .
	\label{eq:potential-U-1st-deriv-space}
\end{equation}
The first-order directional derivative (equation \ref{eq:potential-U-oblique-1st-deriv-space})
coincides with the partial derivatives of first order $\tfrac{\partial}{\partial u} U$ (equation \ref{eq:potential-U-1st-deriv-space})
for the particular case in which the unit vector $\hat{\mathbf{u}}$ coincides with one of the canonical base vectors
\begin{equation}
	\hat{\mathbf{x}} = (1,0,0)^{\top} , \quad 
	\hat{\mathbf{y}} = (0,1,0)^{\top} , \quad
	\hat{\mathbf{z}} = (0,0,1)^{\top} \: .
	\label{eq:canonical-base-vector}
\end{equation}
For example, $D(1, \hat{\mathbf{z}}) \, U(\mathbf{r}) = \tfrac{\partial}{\partial z} U(\mathbf{r})$.

It can also be verified from equation \ref{eq:diff-operators-generic} that 
\begin{equation}
	D(1, \hat{\mathbf{u}}) \, D(1, \hat{\mathbf{v}}) \, U(\mathbf{r}) = 
	\hat{\mathbf{u}}^{\top} \mathbf{H}_{U}(\mathbf{r}) \, \hat{\mathbf{v}} \: ,
	\label{eq:potential-U-oblique-2nd-deriv-space}
\end{equation}
represents a second-order directional derivative along constant unit vectors $\hat{\mathbf{u}}$ and
$\hat{\mathbf{v}}$, where 
\begin{equation}
	\mathbf{H}_{U}(\mathbf{r}) = 
	\begin{bmatrix}
		\tfrac{\partial^{2}}{\partial x^{2}} U(\mathbf{r}) & 
		\tfrac{\partial^{2}}{\partial x \, \partial y} U(\mathbf{r}) & 
		\tfrac{\partial^{2}}{\partial x \, \partial z} U(\mathbf{r}) \\
		\tfrac{\partial^{2}}{\partial x \, \partial y} U(\mathbf{r}) & 
		\tfrac{\partial^{2}}{\partial y^{2}} U(\mathbf{r}) & 
		\tfrac{\partial^{2}}{\partial y \, \partial z} U(\mathbf{r}) \\
		\tfrac{\partial^{2}}{\partial x \, \partial z} U(\mathbf{r}) & 
		\tfrac{\partial^{2}}{\partial y \, \partial z} U(\mathbf{r}) & 
		\tfrac{\partial^{2}}{\partial z^{2}} U(\mathbf{r})
	\end{bmatrix}
	\label{eq:potential-U-Hessian}
\end{equation}
is the Hessian matrix \citep[e.g.,][ p. 158]{kaplan2002} of $U$ (equation \ref{eq:potential-U-space}),
with elements given by the partial derivatives of second order
\begin{equation}
	\tfrac{\partial^{2}}{\partial u \, \partial v} U(\mathbf{r}) = \kappa \iiint\limits_{\mathcal{V}}
	\sigma(\mathbf{r}') \: 
	\tfrac{\partial^{2}}{\partial u \, \partial v} \left( \frac{1}{\| \mathbf{r} - \mathbf{r}'\|} \right)
	\: d\mathcal{V}' \: , \quad \mathbf{r} \notin \overline{\mathcal{V}} \: , \quad u, v \in \{ x, y, z \} \: .
	\label{eq:potential-U-2nd-deriv-space}
\end{equation} 
It can be easily verified that 
the second-order directional derivative $D(1, \hat{\mathbf{u}}) \, D(1, \hat{\mathbf{v}}) \, U(\mathbf{r})$ 
(equation \ref{eq:potential-U-oblique-2nd-deriv-space})
coincides with the partial derivatives of second order $\tfrac{\partial^{2}}{\partial u \, \partial v} U(\mathbf{r})$ 
(equation \ref{eq:potential-U-2nd-deriv-space}) for the particular cases in which the unit vectors 
$\hat{\mathbf{u}}$ and $\hat{\mathbf{v}}$ are given by the canonical base vectors (equation \ref{eq:canonical-base-vector}).
For example, $D(1, \hat{\mathbf{x}}) \, D(1, \hat{\mathbf{y}}) \, U(\mathbf{r}) = \tfrac{\partial^{2}}{\partial x \, \partial y} U(\mathbf{r})$.

\subsection{Some relevant potential fields for geophysical applications}
\label{subsec:relevant-potential-fields}

In gravity exploration methods, the vertical derivative $\tfrac{\partial}{\partial z} U$
(equation \ref{eq:potential-U-1st-deriv-space}) does not define the gravity anomaly, but the \textit{gravity disturbance} data 
\cite[e.g.,][ p. 93]{hofmann-wellenhof-moritz2005}.
The debate around the differences between gravity anomaly and gravity disturbance is beyond the scope of the 
present work and has been carried about from a more geodetic than a geophysical point of view 
\citep[e.g.,][]{li2001, fairhead2003, hackney-featherstone2003, hinze2005, vajda-etal2006, vajda-etal2007, vajda-etal2008}. 
Maybe, that is why it is still unfamiliar to a considerable group of exploration geophysicists.
Another important quantity for gravity exploration method is the Hessian matrix $\mathbf{H}_{U}$
(equation \ref{eq:potential-U-Hessian}), which defines the
(gravitational) \textit{gradient tensor} \citep[e.g.,][]{pedersen_rasmussen1990}.

In magnetic exploration methods, it is well-established that the second-order directional derivative 
$\hat{\mathbf{u}}^{\top} \mathbf{H}_{U} \, \hat{\mathbf{v}}$ (equation \ref{eq:potential-U-oblique-2nd-deriv-space})
is commonly used to approximate the \textit{total-field anomaly} produced by 3D sources \cite[e.g., ][ p. 179]{blakely1996},
where $\hat{\mathbf{u}}$ defines the constant direction of the main geomagnetic
field at the study area and $\hat{\mathbf{v}}$ the constant total-magnetization direction of the sources.
Because of that, we use the term \textit{approximated total-field anomaly} to denote this 
potential field throughout the manuscript.
Note that, in this case, $\hat{\mathbf{u}}^{\top} \mathbf{H}_{U} \, \hat{\mathbf{v}}$ 
(equation \ref{eq:potential-U-oblique-2nd-deriv-space}) can be used to define (i) the  
\textit{Cartesian components of the magnetic induction field} produced due to the 3D
sources by setting the unit vector $\hat{\mathbf{u}}$ equal to the canonical base vectors $\hat{\mathbf{x}}$,
$\hat{\mathbf{y}}$ and $\hat{\mathbf{z}}$ (equation \ref{eq:canonical-base-vector}) or (ii)
the \textit{RTP anomaly} by setting $\hat{\mathbf{u}} = \hat{\mathbf{v}} = \pm\hat{\mathbf{z}}$.
This special anomaly represents the vertical component of 
the magnetic induction field that would be produced by the 3D sources if they had a 
purely induced and vertical magnetization. Differently from generic total-field anomalies 
\cite[e.g.,][ p. 178]{blakely1996}, which have a ``skewed'' pattern 
oscillating from negative to positive values and do not have their maxima placed right 
above the magnetic sources, the RTP anomaly is easier to interpret 
because it is predominantly positive right above the sources and
decays to zero towards their boundaries. \citet{baranov1957} was the first
to propose a method, called \textit{reduction to the pole} (RTP), to transform generic
total-field anomalies into this special anomaly.

\subsection{Classical Dirichlet's and Neumann's problems on a plane}
\label{subsec:dirichlet_neumann}

At all points $\mathbf{r} \notin \overline{\mathcal{V}}$ 
located outside the sources, $U$ (equation
\ref{eq:potential-U-space}) has continuous partial derivatives of all orders, 
satisfies the Laplace's equation and is regular at infinity \citep[][ p. 122--124, 217]{kellogg1967}.
Because of these properties, we have that
\begin{equation}
	U(\mathbf{r}) = 
	\int\limits_{-\infty}^{\infty} \int\limits_{-\infty}^{\infty}
	D(\ell, \hat{\mathbf{z}}) U(\mathbf{r}'') \:\: 
	D(1 - \ell, \hat{\mathbf{z}}) \Psi(\mathbf{r}, \mathbf{r}'')
	\: d\mathcal{S}'' \: , \quad \forall z < z''_{0} < z^{t}_{0} \: ,
	\quad \ell \in \{0, 1\} \: ,
	\label{eq:U-Dirichlet-Neumann}
\end{equation}
where $D(\ell, \hat{\mathbf{u}})$ is the space-domain differential operator 
defined by equation \ref{eq:diff-operators-generic} in terms of the canonical base vector 
$\hat{\mathbf{z}}$ (equation \ref{eq:canonical-base-vector}) and 
\begin{equation}
	\Psi(\mathbf{r}, \mathbf{r}'') = 
	\frac{1}{2\pi} \: 
	\frac{1}{\| \mathbf{r} - \mathbf{r}'' \|} \: , 
	\quad \mathbf{r} \ne \mathbf{r}'' \: ,
	\label{eq:Psi}
\end{equation}
represents the scaled inverse distance function between $\mathbf{r}$ and $\mathbf{r}''$. 
Equation \ref{eq:U-Dirichlet-Neumann} is commonly deduced from 
the \textit{Green's third identity} \citep[][ p. 219]{kellogg1967}, according to
the general outline presented by \cite{roy1962}, for example.

For $\ell = 0$, equation \ref{eq:U-Dirichlet-Neumann} states that $U$ can be defined
in terms of its values $D(0, \hat{\mathbf{u}}) U(\mathbf{r}'') = U(\mathbf{r}'')$ on a horizontal surface $\mathcal{S}''$ 
above the sources and 
represents the solution of the classical \textit{Dirichlet’s problem} or the 
\textit{first boundary value problem of potential theory} \cite[][ p. 236]{kellogg1967} on a plane, where 
$D(1, \hat{\mathbf{z}}) \Psi(\mathbf{r}, \mathbf{r}'') = \tfrac{\partial}{\partial z} \Psi(\mathbf{r}, \mathbf{r}'')$ 
represents the Green's function.
From the physical point of view, equation \ref{eq:U-Dirichlet-Neumann} with $\ell = 0$ defines a continuous 
layer of monopoles (or vertically-magnetized dipoles) that is located on $\mathcal{S}''$, has 
physical-property distribution $U(\mathbf{r}'')$ and exactly reproduces the 
potential field $U(\mathbf{r})$ at any point $\mathbf{r}$ above $\mathcal{S}''$.

On the other hand, equation \ref{eq:U-Dirichlet-Neumann} for $\ell = 1$ shows that $U$ can also
be defined in terms of its vertical derivative 
$D(1, \hat{\mathbf{z}}) U(\mathbf{r}'') = \tfrac{\partial}{\partial z} U(\mathbf{r}'')$ 
a horizontal surface $\mathcal{S}''$ above the sources and represents the solution of the classical \textit{Neumann’s problem} or the 
\textit{second boundary value problem of potential theory} \cite[][ p. 246]{kellogg1967} on a plane, 
with the Green's function given by $D(0, \hat{\mathbf{u}}) \Psi(\mathbf{r}, \mathbf{r}'') = \Psi(\mathbf{r}, \mathbf{r}'')$.
Equation \ref{eq:U-Dirichlet-Neumann} with $\ell = 1$ can be alternatively interpreted as a planar layer of
vertically-magnetized dipoles (or monopoles) that is located on $\mathcal{S}''$, has 
physical-property distribution $\tfrac{\partial}{\partial z} U(\mathbf{r}'')$ and exactly reproduces 
$U(\mathbf{r})$ at any point $\mathbf{r}$ above $\mathcal{S}''$.

By applying the Fourier-convolution theorem \cite[][ p. 268]{blakely1996} to equation \ref{eq:U-Dirichlet-Neumann} 
on a horizontal surface $\mathcal{S}^{\dagger}$ above the sources, we find that
\begin{equation}
	\mathcal{F} [ U(\mathbf{r}^{\dagger}) ] = 
	k \: \mathcal{F} [ U(\mathbf{r}'') ] \:
	\mathcal{F} [ \Psi(\mathbf{r}^{\dagger}, \mathbf{r}'') ] \: ,
	\quad z^{\dagger}_{0} < z''_{0} < z^{t}_{0} \: , \quad k \ne 0 \: ,
	\label{eq:U-Dirichlet-Neumann-Fourier}
\end{equation}
where 
\begin{equation}
	\mathcal{F} [ \Psi(\mathbf{r}^{\dagger}, \mathbf{r}'') ] 
	= \frac{e^{- k \, (z''_{0} - z^{\dagger}_{0})}}{k}
	\: , \quad z^{\dagger}_{0} < z''_{0} \: , \quad k \ne 0 \: ,
	\label{eq:Psi-Fourier}
\end{equation}
is defined according to \citet[][ p. 275]{blakely1996}, with $k = \sqrt{k_{x}^{2} + k_{y}^{2}}$.
We get from equation \ref{eq:diff-operators-generic} that the term $k$ in equation \ref{eq:U-Dirichlet-Neumann-Fourier} represents,
in the Fourier domain, the combined effect of the space-domain differential operators $D(\ell, \hat{\mathbf{z}})$ and 
$D(1-\ell, \hat{\mathbf{z}})$.

\subsection{Potential-field transformations}
\label{subsec:potential_field_transformations}

Several geophysical applications use the observed potential field 
$d^{o}$ (equation \ref{eq:observed-data-space}) to estimate another potential field 
defined in the space domain by
\begin{equation}
	f(\mathbf{r}) = D_{f} \, U(\mathbf{r}) \: , 
	\quad \mathbf{r} \notin \overline{\mathcal{V}} \: ,
	\label{eq:transformed-field-space}
\end{equation}
or, in the Fourier domain, 
\begin{equation}
	\mathcal{F} [ f(\mathbf{r}^{\dagger}) ] =
	\alpha_{f} \, \mathcal{F} [ U(\mathbf{r}^{\dagger}) ] \: ,
	\quad \mathbf{r}^{\dagger} \in \mathcal{S}^{\dagger} \: ,
	\quad z^{\dagger}_{0} < z^{t}_{0} \: ,
	\label{eq:transformed-field-Fourier}
\end{equation}
where $D_{f}$ is a space-domain differential operator acting on $U$ (equation \ref{eq:potential-U-space}) and
$\alpha_{f}$ its corresponding representation in the Fourier domain.
The differential operators $D_{f}$ and $\alpha_{f}$ associated with the transformed potential field 
$f$ (equations \ref{eq:transformed-field-space} and \ref{eq:transformed-field-Fourier}) are commonly
defined in terms of $D$ and $\alpha$ (equation \ref{eq:diff-operators-generic}). 
Estimating $f$ (equations \ref{eq:transformed-field-space} and \ref{eq:transformed-field-Fourier}) from 
the observed potential field $d^{o}$ (equation \ref{eq:observed-data-space}) is an operation 
commonly known in geophysics as \textit{potential-field transformation}.

A common approach to compute potential-field transformations \cite[e.g.,][]{gunn1975} 
consists in using the definition of observed data in the Fourier domain (equation \ref{eq:observed-data-Fourier}) 
and manipulating the Fourier domain solution of the Dirichlet's and Neumann's problems 
(equation \ref{eq:U-Dirichlet-Neumann-Fourier}) to obtain
\begin{equation}
	\mathcal{F} [ f(\mathbf{r}^{\dagger}) ] =
	\left( \tfrac{\alpha_{f}}{\alpha_{d}} \right) \, \mathcal{F} \left[ d^{o}(\mathbf{r}^{\dagger}) \right] \: , 
	\quad \mathbf{r}^{\dagger} \in \mathcal{S}^{\dagger} \: ,
	\quad z^{\dagger}_{0} < z^{t}_{0} \: .
	\label{eq:field-transformation-Fourier}
\end{equation}
Equation \ref{eq:field-transformation-Fourier} defines a single-step potential-field transformation 
in terms of the observed potential field $d^{o}(\mathbf{r}^{\dagger})$ on a planar surface $\mathcal{S}^{\dagger}$ 
located above the sources, at $z^{\dagger}_{0}$. 
The transformed field $f$ in the space domain is obtained by simply computing the inverse Fourier 
transform of \ref{eq:field-transformation-Fourier}.
For the sake of simplicity, we have opted for defining the fields $d^{o}$ and $f$ at the same points in equation 
\ref{eq:field-transformation-Fourier}.
However, it is important to stress that, depending on the 
geophysical application, $f$ may be estimated at points different from those at which the observations are 
placed.

The \textit{equivalent-layer technique} offers an alternative approach that is generally defined in the 
space domain to compute potential-field transformations. What follows aims at discussing the well-established 
empirical approaches used in the equivalent-layer technique, as well as discussing the theoretical background 
supporting them.
\section{Planar equivalent layers}
\label{sec:eqlayers}

The integral equations representing the generalized solutions of Dirichlet's
and Neumann's problems (equations \ref{eq:U-Dirichlet-cross} and \ref{eq:U-Neumann-cross})
can be alternatively interpreted in terms of their kernels.
Note that, according to equation \ref{eq:U-Dirichlet-cross}, the potential 
$U(\mathbf{r})$ is analogous to the vertical component of the gravitational 
attraction produced at all points $\mathbf{r}$ located above $z''_{0}$ by a fictitious 
planar layer of monopoles defined on $z''_{0}$, with mass distribution
$\tau(\mathbf{r}'')$.
This equation also represents a solution of the planar equivalent-layer integral
(equation \ref{eq:integral-equation-eqlayer}) for the particular case in which 
the potential field $d(\mathbf{r})$ is the gravitational potential 
$U(\mathbf{r})$ (equation \ref{eq:scalar-potential-U}) due to the
3D sources and the kernel $a(\mathbf{r}, \mathbf{r}'')$ is defined by 
$\partial_{z} \Psi(\mathbf{r}, \mathbf{r}'')$. In this case, 
$p(\mathbf{r}'') = \tau(\mathbf{r}'')$, which represents (i) the actual gravitational
potential produced by the 3D sources on a plane $z''_{0}$ located above the sources
or (ii) a fictitious gravitational potential obtained by downward continuation to points
inside the sources.

Equation \ref{eq:U-Neumann-cross}, instead, shows that  
$U(\mathbf{r})$ is analogous to the gravitational potential produced at all points
$\mathbf{r}$ located above $z''_{0}$ by a fictitious planar layer of monopoles defined on
$z''_{0}$, with mass distribution $\upsilon(\mathbf{r}'')$.
Besides, it defines another solution of the planar equivalent-layer integral
(equation \ref{eq:integral-equation-eqlayer}) for the case in which the
potential field $d(\mathbf{r})$ is the gravitational potential 
$U(\mathbf{r})$ (equation \ref{eq:scalar-potential-U}) due to the 3D sources
and the kernel is given by $a(\mathbf{r}, \mathbf{r}'') = \Psi(\mathbf{r}, \mathbf{r}'')$.
In this case, the continuous physical property distribution $p(\mathbf{r}'')$ is defined
by $\upsilon(\mathbf{r}'')$, which represents (i) the actual vertical component of the
gravitational attraction produced by the 3D sources on a 
plane $z''_{0}$ located above the sources or (ii) a fictitious gravity disturbance 
obtained by downward continuation to points inside the sources.


\subsection{Planar equivalent layer of monopoles applied to gravity data}
\label{subsec:eqlayer_monopoles}


The integral equations representing the generalized solutions of Dirichlet's
and Neumann's problems (equations \ref{eq:U-Dirichlet-cross} and \ref{eq:U-Neumann-cross})
represent solutions for the planar equivalent-layer integral for the case in which
the potential field $d(\mathbf{r})$ is the gravitational potential 
$U(\mathbf{r})$ (equation \ref{eq:scalar-potential-U}) due to the 3D sources.
In practical situation, however, the potential-field measurements are represented
by the partial derivatives of first and second orders of $U(\mathbf{r})$
(equations \ref{eq:1st-deriv-U} and \ref{eq:2nd-deriv-U}).
Here, we show the planar equivalent layers associated with these derivatives can be
obtained from the generalized solution of Neumann's problem 
(equation \ref{eq:U-Neumann-cross}).

Consider the partial derivatives of first and second orders of $U(\mathbf{r})$ obtained
by differentiating equation \ref{eq:U-Neumann-cross}:
\begin{equation}
\partial_{\alpha} U(\mathbf{r}) =  
\int\limits_{-\infty}^{\infty} \int\limits_{-\infty}^{\infty}
\upsilon(\mathbf{r}'') \:\: 
\partial_{\alpha} \Psi(\mathbf{r}, \mathbf{r}'') \: dx''dy'' \: , 
\quad z < z''_{0} \: ,
\quad \alpha = x, y, z \: ,
\label{eq:1st-deriv-U-Neumann}
\end{equation}
and
\begin{equation}
\partial_{\alpha\beta} U(\mathbf{r}) =  
\int\limits_{-\infty}^{\infty} \int\limits_{-\infty}^{\infty}
\upsilon(\mathbf{r}'') \:\: 
\partial_{\alpha\beta} \Psi(\mathbf{r}, \mathbf{r}'') \: dx''dy'' \: , 
\quad z < z''_{0} \: ,
\quad \alpha, \beta = x, y, z \: .
\label{eq:2nd-deriv-U-Neumann}
\end{equation}
Note that the derivatives are computed with respect to the coordinates
of $\mathbf{r}$ and, consequently, act only on the kernel $\Psi(\mathbf{r}, \mathbf{r}'')$
and not on the function $\upsilon(\mathbf{r}'')$ 
(equation \ref{eq:U-Neumann-cross}).

Note that equation \ref{eq:1st-deriv-U-Neumann} represents a solution of the
planar equivalent layer integral for the case in which the potential field 
$d(\mathbf{r})$ is the vertical component $\partial_{z} U(\mathbf{r})$ of the
gravitational attraction due to the 3D sources, which is commonly used to
approximate gravity disturbance data \cite[e.g.,][ p. 93]{hofmann-wellenhof-moritz2005}.
In this case, the kernel and physical property distribution are given by 
$a(\mathbf{r}, \mathbf{r}'') = \partial_{z} \Psi(\mathbf{r}, \mathbf{r}'')$ and
$p(\mathbf{r}'') = \upsilon(\mathbf{r}'')$, respectively.
This result shows that there exists a continuous planar
equivalent layer of monopoles that exactly reproduces the gravity disturbance
data above the sources.
\citet{dampney1969} was the pioneer in using this result to interpolate gravity data and
a considerable amount of literature has been published on this topic since then.
This solution $p(\mathbf{r}'')$ can also be directly deduced from the
\textit{upward-continuation integral}
\citep[e.g., ][]{peters1949, roy1962, bhattacharyya1967, henderson1970, gunn1975}, 
which represents the solution of Dirichlet’s problem on a plane
(equation \ref{eq:U-Dirichlet}), or from the generalized solution presented here
(equation \ref{eq:U-Dirichlet}-cross).

Consider now the case in which the potential-field data $d(\mathbf{r})$ in the 
planar equivalent layer integral (equation \ref{eq:integral-equation-eqlayer}) 
are gravitational gradient tensor components. These data are usually described as the
partial derivatives of second order of the gravitational potential. Hence, by using the
second derivatives defined by equation \ref{eq:2nd-deriv-U-Neumann}, we obtain for the 
$\alpha\beta$-component of gravity tensor that
$d(\mathbf{r}) = \partial_{\alpha\beta} U(\mathbf{r})$ and 
$a(\mathbf{r}, \mathbf{r}'') = \partial_{\alpha\beta} \Psi(\mathbf{r}, \mathbf{r}'')$, 
so that the continuous equivalent layer has a closed-form solution 
$p(\mathbf{r}'') = \upsilon(\mathbf{r}'')$.
This result shows that the same continuous planar equivalent layer of monopoles that 
reproduces the observed gravity disturbance data also reproduces the 
gradient tensor data above the sources.
This result has been used by the scientific community for processing gravity-gradient
data \citep{barnes-lumley_2011} or converting gravity disturbance into the
gradient tensor \citep{piauilino_etal2019_seg, li_etal2021}, for example.
However, researchers have not treated this topic in much detail and it is usually assumed
to be true without any proof.

\subsection{Planar equivalent layer of dipoles applied to magnetic data}
\label{subsec:eqlayer_dipoles}

We may also combine the partial derivatives of first order 
(equation \ref{eq:1st-deriv-U-Neumann}) to compute the
following directional derivative of first order:
\begin{equation}
\partial_{h} U(\mathbf{r}) = 
\int\limits_{-\infty}^{\infty} \int\limits_{-\infty}^{\infty}
\upsilon(\mathbf{r}'') \:\: 
\partial_{h} \Psi(\mathbf{r}, \mathbf{r}'') \:\: dx''dy'' \: , 
\quad z < z''_{0} \: ,
\quad \alpha, \beta = x, y, z \: ,
\label{eq:oblique-1st-deriv-U-Neumann}
\end{equation}
where 
\begin{equation}
\partial_{h} \Psi(\mathbf{r}, \mathbf{r}'') = 
\nabla \Psi(\mathbf{r}, \mathbf{r}'')^{\top} \hat{\mathbf{h}} 
\label{eq:oblique-1st-deriv-Psi}
\end{equation}
is the directional derivative of the inverse distance function 
$\Psi(\mathbf{r}, \mathbf{r}'')$ (equation \ref{eq:Psi}) along the
unit vector $\hat{\mathbf{h}}$ and
$\nabla \Psi(\mathbf{r}, \mathbf{r}'')$ is the gradient of $\Psi(\mathbf{r}, \mathbf{r}'')$
with respect to the coordinates of the observation point $\mathbf{r}$.

Finally, we may use partial derivatives of second order 
(equation \ref{eq:2nd-deriv-U-Neumann}) to obtain the following directional derivative
of second order:
\begin{equation}
\partial_{th} U(\mathbf{r}) = 
\int\limits_{-\infty}^{\infty} \int\limits_{-\infty}^{\infty}
\upsilon(\mathbf{r}'') \:\: 
\hat{\mathbf{t}}^{\top} \mathbf{H}_{\Psi}(\mathbf{r}, \mathbf{r}'') \hat{\mathbf{h}} 
\:\: dx''dy'' \: , \quad z < z''_{0} \: ,
\quad \alpha, \beta = x, y, z \: ,
\label{eq:oblique-2nd-deriv-Theta-Neumann}
\end{equation}
where
\begin{equation}
\mathbf{H}_{\Psi}(\mathbf{r}, \mathbf{r}'') = 
\begin{bmatrix}
\partial_{xx} \Psi(\mathbf{r}, \mathbf{r}'') & 
\partial_{xy} \Psi(\mathbf{r}, \mathbf{r}'') & 
\partial_{xz} \Psi(\mathbf{r}, \mathbf{r}'') \\
\partial_{xy} \Psi(\mathbf{r}, \mathbf{r}'') & 
\partial_{yy} \Psi(\mathbf{r}, \mathbf{r}'') & 
\partial_{yz} \Psi(\mathbf{r}, \mathbf{r}'') \\
\partial_{xz} \Psi(\mathbf{r}, \mathbf{r}'') & 
\partial_{yz} \Psi(\mathbf{r}, \mathbf{r}'') & 
\partial_{zz} \Psi(\mathbf{r}, \mathbf{r}'')
\end{bmatrix} \: ,
\label{eq:H-Psi}
\end{equation}
is the Hessian matrix of the inverse distance function 
$\Psi(\mathbf{r}, \mathbf{r}'')$ (equation \ref{eq:Psi}) defined in terms of
the second derivatives $\partial_{\alpha\beta} \Psi(\mathbf{r}, \mathbf{r}'') \equiv 
\frac{\partial^{2}}{\partial\alpha \, \partial\beta} \Psi(\mathbf{r}, \mathbf{r}'')$,
$\alpha = x, y, z$, $\beta = x, y, z$, with respect to the coordinates of
the observation point $\mathbf{r}$.
\section{Existence of continuous equivalent layers on a plane}
\label{sec:planar_eqlayer_theorem}

\subsection{Modified Dirichlet's and Neumann's problems on a plane}

To prove the existence of continuous equivalent layers $p_{\varepsilon}(\mathbf{r}'')$
according to Definition \ref{def:continuous-eqlayer},
we need first modify the Dirichlet’s and Neumann’s problems 
(equations \ref{eq:U-Dirichlet-Neumann} and \ref{eq:U-Dirichlet-Neumann-Fourier}) for the case in which
the plane $z''_{0}$ cross the sources.
Hence, let us first consider an approximating potential $U_{\varepsilon}(\mathbf{r}'')$ 
defined on a horizontal plane $z''_{0}$ that may be located above or cross 
the sources (Figure \ref{fig:geometry}), with Fourier transform given by
\begin{equation}
	\mathcal{F} \left[ U_{\varepsilon}(\mathbf{r}'') \right] = 
	\begin{cases}
		\mathcal{F} \left[ U(\mathbf{r}^{\dagger}) \right] 
		e^{\, k \, (z''_{0} - z^{\dagger}_{0})} & , 
		\quad 0 < k \le k_{\varepsilon} \\
		0 & , \quad k > k_{\varepsilon} 
	\end{cases} \: ,
	\label{eq:potential-U-epsilon-Fourier-layer}
\end{equation}
where $k_{\varepsilon}$ is a fixed and possibly very large truncating wavenumber and $z^{\dagger}_{0} < z^{t}_{0}$. 
Consider also that the inverse Fourier transform of $\mathcal{F} \left[ U_{\varepsilon}(\mathbf{r}'') \right]$ 
(equation \ref{eq:potential-U-epsilon-Fourier-layer}) satisfies
\begin{equation}
	U_{\varepsilon}(\mathbf{r}) = 
	\int\limits_{-\infty}^{\infty} \int\limits_{-\infty}^{\infty}
	D(\ell, \hat{\mathbf{z}}) U_{\varepsilon}(\mathbf{r}'') \:\:
	D(1-\ell, \hat{\mathbf{z}}) \Psi(\mathbf{r}, \mathbf{r}'') 
	\:\: d\mathcal{S}'' \: , 
	\quad \forall z < z^{t}_{0} \: , \quad \forall z < z''_{0} \: ,
	\quad \ell \in \{ 0, 1 \} \: ,
	\label{eq:potential-U-epsilon-space}
\end{equation}
where $D$ is the space-domain differential operator 
defined by equation \ref{eq:diff-operators-generic} in terms of the canonical base vector 
$\hat{\mathbf{z}}$ (equation \ref{eq:canonical-base-vector}).

Note that the function $U_{\varepsilon}(\mathbf{r}'')$ obtained by computing
the inverse Fourier transform of 
$\mathcal{F} \left[ U_{\varepsilon}(\mathbf{r}'') \right]$ (equation \ref{eq:potential-U-epsilon-Fourier-layer})
approximates the potential field obtained by 
computing the downward continuation of $U$ (equation \ref{eq:potential-U-space}) 
to points that may be located inside the sources. 
It is well-known, however, that downward continuation can only be computed at 
source-free points, where the Laplace's equation applies. This means that the
continuation plane is limited by the top $z^{t}_{0}$ of the sources, otherwise the continued field 
presents an oscillatory character \citep[e.g.,][]{roy1967, fedi_florio2002} that may
prevent $\mathcal{F} \left[ U_{\varepsilon}(\mathbf{r}'') \right]$ to be square integrable.
The truncating wavenumber $k_{\varepsilon}$ (equation \ref{eq:potential-U-epsilon-Fourier-layer}) 
overcomes this problem, as suggested by \citet{bhattacharyya1967},
and guarantees that the Fourier-convolution theorem 
can be applied to the auxiliary potential $U_{\varepsilon}$
(equation \ref{eq:potential-U-epsilon-space}) on a horizontal plane $z^{\dagger}_{0} < z^{t}_{0}$ 
to obtain
\begin{equation}
	\mathcal{F} [ U_{\varepsilon}(\mathbf{r}^{\dagger}) ]
	= k \, \mathcal{F} [ U_{\varepsilon}(\mathbf{r}'') ] \:
	\mathcal{F} [ \Psi(\mathbf{r}^{\dagger}, \mathbf{r}'') ] \: .
	\label{eq:potential-U-epsilon-Fourier-dagger}
\end{equation}
We get by substituting \ref{eq:Psi-Fourier} and \ref{eq:potential-U-epsilon-Fourier-layer} into 
\ref{eq:potential-U-epsilon-Fourier-dagger} that
\begin{equation}
	\mathcal{F} [ U_{\varepsilon}(\mathbf{r}^{\dagger}) ] = 
	\begin{cases}
		\mathcal{F} [ U(\mathbf{r}^{\dagger}) ] &
		, \quad 0 < k \le k_{\varepsilon} \\
		0 & , \quad k > k_{\varepsilon} 
	\end{cases} \: ,
	\label{eq:potential-U-epsilon-Fourier-dagger2}
\end{equation}
which shows that the Fourier transforms of 
$U(\mathbf{r}^{\dagger})$ (equation \ref{eq:potential-U-Fourier}) and 
$U_{\varepsilon}(\mathbf{r}^{\dagger})$ (equation \ref{eq:potential-U-epsilon-space}) 
on the horizontal plane $z^{\dagger}_{0}$ located above the 3D sources coincide up 
to the truncating wavenumber $k_{\varepsilon}$.

We have just specified the Dirichlet's and Neumann's problems 
(equations \ref{eq:U-Dirichlet-Neumann} and \ref{eq:U-Dirichlet-Neumann-Fourier}) in terms of a truncating wavenumber 
$k_{\varepsilon}$ with the purpose of expanding their applicability to the case in which plane $z''_{0}$ cross the sources. 
In this case, we obtain an approximating potential $U_{\varepsilon}(\mathbf{r})$ that coincides with the true 
potential $U(\mathbf{r})$ up to the truncating wavenumber $k_{\varepsilon}$ 
(equation \ref{eq:potential-U-epsilon-Fourier-dagger2}). 
At last, it can be easily verified that, if we let the truncating wavenumber $k_{\varepsilon}$ 
goes to infinite in equation \ref{eq:potential-U-epsilon-Fourier-layer} for the case in which $z''_{0} < z^{t}_{0}$, we get 
\begin{equation}
	\lim_{k_{\varepsilon}\to\infty} \mathcal{F} \left[ U_{\varepsilon}(\mathbf{r}'') \right] = 
	\mathcal{F} \left[ U(\mathbf{r}'') \right] \: , \quad z''_{0} < z^{t}_{0} \: ,
	\label{eq:lim-kepsilon-potential-U-epsilon-Fourier-layer}
\end{equation}
\begin{equation}
	\lim_{k_{\varepsilon}\to\infty} \mathcal{F} \left[ U_{\varepsilon}(\mathbf{r}^{\dagger}) \right] = 
	\mathcal{F} \left[ U(\mathbf{r}^{\dagger}) \right] \: , \quad z^{\dagger}_{0} < z''_{0} < z^{t}_{0} \: ,
	\label{eq:lim-kepsilon-potential-U-epsilon-Fourier-dagger}
\end{equation}
pointwise and, consequently, that
\begin{equation}
	\lim_{k_{\varepsilon}\to\infty} U_{\varepsilon}(\mathbf{r}) = 
	U(\mathbf{r}) \: , \quad \forall z < z''_{0} < z^{t}_{0} \: ,
	\label{eq:lim-kepsilon-potential-U-epsilon-space}
\end{equation}
which is in perfect agreement with the solutions of the Dirichlet's and Neumann's
problems (equations \ref{eq:U-Dirichlet-Neumann} and \ref{eq:U-Dirichlet-Neumann-Fourier}) on a plane $z''_{0}$ 
that is located above the sources.

\subsection{The planar equivalent-layer theorem}

\begin{theorem}[\textbf{Planar equivalent-layer theorem}]
	\label{theorem:planar-eqlayer-theorem}
	Let a planar surface $\mathcal{S}''$ located at $z''_{0}$ and a kernel function $a(\mathbf{r}, \mathbf{r}'')$ 
	be set according to Definition \ref{def:continuous-eqlayer}, for a given observed potential field $d^{o}$. 
	Then, for every $\varepsilon > 0$, there exists a continuous equivalent layer $p_{\varepsilon}$. Furthermore, if 
	$z''_{0} < z^{t}_{0}$, then the continuous equivalent layer reproduces $d^{o}$ exactly.
\end{theorem}
\begin{proof}
	We start by defining $p_{\varepsilon}$ (Definition \ref{def:continuous-eqlayer}) 
	in the Fourier domain via
	\begin{equation}
		\mathcal{F} [ p_{\varepsilon}(\mathbf{r}'') ] =
		\tfrac{k \, \alpha_{d}}{\alpha_{a}} \,
		\mathcal{F} [ U_{\varepsilon}(\mathbf{r}'') ] \: ,
		\label{eq:p-Fourier}
	\end{equation}
	where $U_{\varepsilon}(\mathbf{r}'')$ is defined as in \ref{eq:potential-U-epsilon-Fourier-layer}, 
	for a given truncating wavenumber $k_{\varepsilon}$, and $\alpha_{a}$ is differential operator 
	defined in the Fourier domain in terms of $\alpha$ (equation \ref{eq:diff-operators-generic}).
	For any finite choice of $0 < k_{\varepsilon} < \infty$, both $\mathcal{F} [ p_{\varepsilon}(\mathbf{r}'') ]$ 
	and $p_{\varepsilon}(\mathbf{r}'')$ are square-integrable and, consequently, it is possible to
	apply the Fourier-convolution theorem to \ref{eq:predicted-data-continuous-layer} and obtain
	\begin{equation}
		\mathcal{F} [ d_{\varepsilon}(\mathbf{r}^{\dagger}) ] =
		\mathcal{F} [ p_{\varepsilon}(\mathbf{r}'') ] \, 
		\mathcal{F} [ a(\mathbf{r}^{\dagger}, \mathbf{r}'') ] \: ,
		\quad z^{\dagger}_{0} < z''_{0} \: ,
		\label{eq:predicted-data-continuous-layer-Fourier}
	\end{equation}
	where 
	\begin{equation}
		\mathcal{F} [ a(\mathbf{r}^{\dagger}, \mathbf{r}'') ] = \alpha_{a} \, \mathcal{F} [ \Psi(\mathbf{r}^{\dagger}, \mathbf{r}'') ] \: ,
		\label{eq:kernel-a-Fourier}
	\end{equation}
	with $\mathbf{r}^{\dagger} = (x^{\dagger}, y^{\dagger}, z^{\dagger}_{0})^{\top} \notin \overline{\mathcal{V}}$ 
	defining a point on the horizontal surface $\mathcal{S}^{\dagger}$ above the sources, at $z^{\dagger}_{0} < z''_{0}$.
	By considering that the kernel function is given by $a(\mathbf{r}, \mathbf{r}'') = D_{a} \Psi(\mathbf{r}, \mathbf{r}'')$
	and using \ref{eq:observed-data-Fourier} and \ref{eq:Psi-Fourier}, we find that $d_{\varepsilon}$ satisfies
	\begin{equation}
		\mathcal{F} [ d_{\varepsilon}(\mathbf{r}^{\dagger}) ] = 
		\begin{cases}
			\mathcal{F} [ d^{o}(\mathbf{r}^{\dagger}) ] &
			, \quad 0 < k \le k_{\varepsilon} \\
			0 & , \quad k > k_{\varepsilon} 
		\end{cases} \: .
		\label{eq:predicted-data-continuous-layer-Fourier2}
	\end{equation}
	From Plancherel's theorem \citep[][ p. 119]{bracewell1999} and use of polar coordinates 
	$k_{x} = k \, \cos\theta$, $k_{y} = k \, \sin\theta$ and $k = \sqrt{k_{x}^{2} + k_{y}^{2}}$,
	we get
	\begin{equation}
		\begin{split}
		\int\limits_{-\infty}^{\infty} \int\limits_{-\infty}^{\infty}
		\vert 
		d^{o}(\mathbf{r}^{\dagger}) - d(\mathbf{r}^{\dagger})
		\vert^{2} 
		\: d\mathcal{S}^{\dagger} = 
		\int\limits_{-\infty}^{\infty} \int\limits_{-\infty}^{\infty}
		\vert \mathcal{F} [ d^{o}(\mathbf{r}^{\dagger}) ] - \mathcal{F} [
		d(\mathbf{r}^{\dagger}) ] \vert^{2} 
		\: d k_{x} \, d k_{y} \\
		= \int\limits_{0}^{2\pi} \int\limits_{0}^{\infty}
		\vert \mathcal{F} [ d^{o}(\mathbf{r}^{\dagger}) ] - \mathcal{F} [
		d(\mathbf{r}^{\dagger}) ] \vert^{2} \: k \, dkd\theta =
		\int\limits_{0}^{2\pi} \int\limits_{k_{\varepsilon}}^{\infty}
		\vert \mathcal{F} [ d^{o}(\mathbf{r}^{\dagger}) ] \vert^{2} 
		\: k \: d k \, d \theta \: .
		\end{split}
		\label{eq:Plancherel}
	\end{equation}
	Since $\mathcal{F} [ d^{o}(\mathbf{r}^{\dagger}) ] \, k$ is square integrable on $\mathcal{S}^{\dagger}$, we may choose 
	a finite but sufficiently large $k_{\varepsilon} < \infty$ such that 
	$\int_{-\infty}^{\infty} \int_{-\infty}^{\infty} \vert d^{o}(\mathbf{r}^{\dagger}) - d(\mathbf{r}^{\dagger}) 
	\vert^{2} \: d\mathcal{S}^{\dagger} = \int_{0}^{2\pi} \int_{0}^{\infty} 
	\vert \mathcal{F} [ d^{o}(\mathbf{r}^{\dagger}) ] - \mathcal{F} [d(\mathbf{r}^{\dagger}) ] \vert^{2} \: k \, dkd\theta 
	\le \varepsilon$. This leads to the conclusion that $p_{\varepsilon}$ is a continuous equivalent layer. 
	If the layer plane is above the sources (i.e., if $z''_{0} < z^{t}_{0}$), then both 
	$\mathcal{F} [ p_{\varepsilon}(\mathbf{r}'') ]$ and $p_{\varepsilon}(\mathbf{r}'')$ are square-integrable
	even for the choice $k_{\varepsilon} = \infty$. Then, in that case, 
	$\int_{-\infty}^{\infty} \int_{-\infty}^{\infty} \vert d^{o}(\mathbf{r}^{\dagger}) - d(\mathbf{r}^{\dagger}) 
	\vert^{2} \: d\mathcal{S}^{\dagger} = \int_{0}^{2\pi} \int_{0}^{\infty} 
	\vert \mathcal{F} [ d^{o}(\mathbf{r}^{\dagger}) ] - \mathcal{F} [d(\mathbf{r}^{\dagger}) ] \vert^{2} \: k \, dkd\theta 
	= 0$ and $p_{\varepsilon}$ yields a predicted field $d_{\varepsilon}$ that exactly reproduces the observed field $d^{o}$ 
	on $\mathcal{S}^{\dagger}$.
\end{proof}
	
\subsubsection{Remark 1}

The above theorem guarantees the existence of an approximate equivalent layer, even if
the layer crosses the sources, i.e., if $z''_{0} > z^{t}_{0}$.
But it does not provide any guidance on the particular choice of the truncating wave number $k_{\varepsilon}$ 
for the construction of $p_{\varepsilon}$.
However, it is possible to combine the upper bounds defined in Appendix A 
with the right-hand side of equation \ref{eq:Plancherel} to show that
\begin{equation}
	\int\limits_{0}^{2\pi} \int\limits_{0}^{\infty} 
	\vert \mathcal{F} [ d^{o}(\mathbf{r}^{\dagger}) ] - \mathcal{F} [d(\mathbf{r}^{\dagger}) ] \vert^{2} \: k \, dkd\theta 
	\le \varepsilon(q, k_{\varepsilon}, z^{\dagger}_{0}) \: ,
	\label{eq:Plancherel-epsilon}
\end{equation}
where
\begin{equation}
	\varepsilon(q, k_{\varepsilon}, z^{\dagger}_{0}) = \tfrac{2\pi \, \Omega^{2} \, 3^{q}}{(2 \, \Delta z)^{2q}} \: 
	\Gamma(2 \, q \: , \: 2 \, \Delta z \, k_{\varepsilon}) \: ,
	\quad \Delta z = z^{t}_{0} - z^{\dagger}_{0} > 0 \: ,
	\label{eq:epsilon}
\end{equation}
is a real function depending on $\Omega$ (Apendix A, equation \ref{eq:Omega}) and the Incomplete Gamma function 
\citep[][ p. 899]{gradshteyn_ryzhik2007}
\begin{equation}
	\Gamma(2 \, q \: , \: 2 \, \Delta z \, k_{\varepsilon}) = 
	\int\limits_{k_{\varepsilon}}^{\infty}
	e^{- 2 \, \Delta z \, k } \, k^{2q-1} \, (2 \, \Delta z)^{2q} 
	\: d k \: ,
	\label{eq:incomplete-Gamma-function}
\end{equation}
with $q$ being the order of the differential operator $\alpha_{d}$ (equation \ref{eq:diff-operators-generic}). 
According to equation \ref{eq:epsilon}, if we have some a priori bounds 
on the physical-property distribution and position of the true sources, we get a connection between the choice of 
$k_{\varepsilon}$ and the desired accuracy level $\varepsilon$. 
This can potentially support an estimation of the truncating wave number $k_{\varepsilon}$.

\subsubsection{Remark 2}

The truncating wavenumber $k_{\varepsilon}$ plays a central role in the planar equivalent-layer theorem 
(Theorem \ref{theorem:planar-eqlayer-theorem}).
We would like to highlight here that the interpreter typically does not explicitly set $k_{\varepsilon}$ in
practical situations. In fact, obtaining a discrete equivalent layer (Definition \ref{def:discrete-eqlayer}) 
does not strictly imply a truncation of the wavenumber at some maximal $k_{\varepsilon}$, 
but it leads to an exponentially fast decay at sufficiently large wavenumbers. 
From a theoretical point of view, it might be interesting to study the influence of
such a discretization and its relation to $k_{\varepsilon}$ in more detail. 
For practical purposes it seems that this discretization is the predominant factor rendering 
the equivalent-layer technique applicably even when the layer plane crosses the sources.

\subsubsection{Remark 3}

The equivalent layer from Theorem \ref{theorem:planar-eqlayer-theorem} leads to a predicted 
potential field that does not satisfy
the error bound \ref{eq:epsilon-bound-data} only on the surface $\mathcal{S}^{\dagger}$, but also 
on any $\mathcal{S}^{\ddagger}$ with $z^{\ddagger}_{0} < z^{\dagger}_{0}$. 
Namely, it follows from equation \ref{eq:potential-U-Fourier} that
\begin{equation}
	\int\limits_{-\infty}^{\infty} \int\limits_{-\infty}^{\infty}
	\vert d^{o}(\mathbf{r}^{\ddagger}) - d(\mathbf{r}^{\ddagger})
	\vert^{2} \: d\mathcal{S}^{\ddagger} \le
	\int\limits_{-\infty}^{\infty} \int\limits_{-\infty}^{\infty}
	\vert d^{o}(\mathbf{r}^{\dagger}) - d(\mathbf{r}^{\dagger})
	\vert^{2} \: d\mathcal{S}^{\dagger} \: ,
	\quad z^{\ddagger}_{0} < z^{\dagger}_{0} \: .
	\label{eq:error-bound-upper-planes}
\end{equation}
Thus, if $p_{\varepsilon}$ produces a $d_{\varepsilon}$ satisfying 
$\int_{-\infty}^{\infty} \int_{-\infty}^{\infty} \vert d^{o}(\mathbf{r}^{\dagger}) - d(\mathbf{r}^{\dagger}) 
\vert^{2} \: d\mathcal{S}^{\dagger} \le \varepsilon$, then it automatically holds that 
$\int_{-\infty}^{\infty} \int_{-\infty}^{\infty} \vert d^{o}(\mathbf{r}^{\ddagger}) - d(\mathbf{r}^{\ddagger}) 
\vert^{2} \: d\mathcal{S}^{\ddagger} \le \varepsilon$.

\section{Potential-field transformations via equivalent-layer technique}

Note that, given a particular equivalent layer $\tilde{p}_{\varepsilon}$ (Definition \ref{def:continuous-eqlayer}) 
defined in terms of specific $\alpha_{d}$ and $\alpha_{a}$ as a square-integrable function by equation \ref{eq:p-Fourier}, 
it is theoretically possible 
to compute any potential field transformation (equation \ref{eq:transformed-field-space} and \ref{eq:transformed-field-Fourier}).
This can be verified by first noting that applying the Fourier-convolution theorem to \ref{eq:transformed-field-continuous-layer} 
on a planar surface $\mathcal{S}^{\ddagger}$ at a constant vertical coordinate $z^{\ddagger}_{0} < z''_{0}$ gives
\begin{equation}
	\mathcal{F} [ f_{\varepsilon}(\mathbf{r}^{\ddagger}) ] =
	\mathcal{F} \left[ \tilde{p}_{\varepsilon}(\mathbf{r}'') \right]
	\mathcal{F} [ t(\mathbf{r}^{\ddagger}, \mathbf{r}'') ] \: , 
	\quad z^{\ddagger}_{0} < z''_{0} \: ,
	\label{eq:transformed-field-continuous-layer-Fourier}
\end{equation}
where
\begin{equation}
	\mathcal{F} [ t(\mathbf{r}^{\ddagger}, \mathbf{r}'') ] = \alpha_{t} \, \mathcal{F} [ \Psi(\mathbf{r}^{\ddagger}, \mathbf{r}'') ]
	\label{eq:kernel-t-Fourier}
\end{equation}
is defined in terms of a differential operator $\alpha_{t}$ acting on $\Psi$ (equation \ref{eq:Psi}) in the Fourier domain.
Then we set 
\begin{equation}
	\alpha_{t} = \left( \tfrac{\alpha_{a}}{\alpha_{d}} \right) \alpha_{f} \: ,
	\label{eq:alpha-t-generic}
\end{equation}
where $\alpha_{f}$ is the differential operator associated with the desired transformed field $f$ (equation \ref{eq:transformed-field-Fourier}).
In this case, we get from equations \ref{eq:Psi-Fourier}, \ref{eq:p-Fourier}  and \ref{eq:kernel-t-Fourier} 
that the transformed potential field $f_{\varepsilon}$ (equation \ref{eq:transformed-field-continuous-layer-Fourier}) 
is given by:
\begin{equation}
	\mathcal{F} [ f_{\varepsilon}(\mathbf{r}^{\ddagger}) ] =
	\alpha_{f} \, k \, \mathcal{F} \left[ U_{\varepsilon}(\mathbf{r}'') \right]
	\mathcal{F} [ \Psi(\mathbf{r}^{\ddagger}, \mathbf{r}'') ] \: , 
	\quad z^{\ddagger}_{0} < z''_{0} \: .
	\label{eq:transformed-field-continuous-layer-Fourier2}
\end{equation}
Finally, we use \ref{eq:Psi-Fourier}, \ref{eq:transformed-field-Fourier} and \ref{eq:potential-U-epsilon-Fourier-layer} to show that
\begin{equation}
	\mathcal{F} [ f_{\varepsilon}(\mathbf{r}^{\ddagger}) ] = 
	\begin{cases}
		\mathcal{F} [ f(\mathbf{r}^{\ddagger}) ] &
		, \quad 0 < k \le k_{\varepsilon} \\
		0 & , \quad k > k_{\varepsilon} 
	\end{cases} \: ,
	\label{eq:transformed-field-continuous-layer-Fourier3}
\end{equation}
which shows that the transformed field $f_{\varepsilon}$ yielded by the equivalent layer coincides with 
that produced by the true sources on the planar surface $\mathcal{S}^{\ddagger}$ up 
to the truncating wavenumber $k_{\varepsilon}$.

In the context of the equivalent-layer technique, a potential-field transformation is generally defined
in the space domain, by following two steps. The first consists in obtaining a particular layer $\tilde{p}_{\varepsilon}$ 
(Definition \ref{def:continuous-eqlayer}), which is subsequently used to compute a transformed field $f_{\varepsilon}$ via 
\ref{eq:transformed-field-continuous-layer}.
It can be verified, however, that this two-step approach is equivalent to the single-step approach based on the 
Fourier transform of the observed potential field (equation \ref{eq:field-transformation-Fourier}). 
We find by using \ref{eq:observed-data-Fourier} to manipulate \ref{eq:transformed-field-continuous-layer-Fourier2} that
\begin{equation}
	\mathcal{F} [ f_{\varepsilon}(\mathbf{r}^{\ddagger}) ] =
	\left( \tfrac{\alpha_{f}}{\alpha_{d}} \right) \, 
	\begin{cases}
		\mathcal{F} \left[ d^{o}(\mathbf{r}^{\dagger}) \right] 
		e^{\, k \, (z^{\ddagger}_{0} - z^{\dagger}_{0})} & , \quad 0 < k \le k_{\varepsilon} \\
		0 & , \quad k > k_{\varepsilon} 
	\end{cases} \: .
	\label{eq:transformed-field-continuous-layer-single-step}
\end{equation}
Note that, for the particular case in which $z^{\ddagger}_{0}$ coincides with $z^{\dagger}_{0}$ and 
the truncating wavenumber $k_{\varepsilon}$ goes to infinite, equation 
\ref{eq:transformed-field-continuous-layer-single-step} reduces to \ref{eq:field-transformation-Fourier}.
This result establishes a theoretical link between potential-field transformations made via
equivalent-layer technique and those commonly made via Fourier transform of the observed data.
This link explains, for example, why potential-field transformations made via equivalent-layer technique may suffer from the same 
numerical problems generally associated with those made via Fourier transform of the observed potential-field data, 
as have already been pointed out by \citet{li_etal2014} about the reduction to the pole.

\section{Equivalent layer classes}
\label{sec:eqlayer_classes}

\subsection{Dirichlet's equivalent layers}

\begin{definition}
	An equivalent layer is called Dirichlet's equivalent layer if the Fourier transform of its physical-property distribution has 
	the form
	\begin{equation}
		\mathcal{F} [ p(\mathbf{r}'') ] =
		\begin{cases}
			\alpha_{d} \, \mathcal{F} [ U(\mathbf{r}^{\dagger}) ] 
			e^{\, k \, (z''_{0} - z^{\dagger}_{0})} & , 
			\quad 0 < k \le k_{\varepsilon} \\
			0 & , \quad k > k_{\varepsilon} 
		\end{cases} \: .
		\label{eq:p-Fourier-Dirichlet}
	\end{equation}
	where $z''_{0} > z^{\dagger}_{0}$.
	\label{def:eql-Dirichlet}
\end{definition}

Note that the particular $\mathcal{F} [ p(\mathbf{r}'') ]$ associated with Dirichlet's layers (equation \ref{eq:p-Fourier-Dirichlet}) 
is defined from the general solution (equation \ref{eq:p-Fourier}) by setting $\alpha_{a} \equiv \alpha(1, \hat{\mathbf{z}})$ 
(equations \ref{eq:diff-operators-generic}, \ref{eq:canonical-base-vector} and \ref{eq:kernel-Fourier}).
It is important noting that this particular $\mathcal{F} [ p(\mathbf{r}'') ]$ depends on $\alpha_{d}$, which 
carries information about the observed potential field (equation \ref{eq:observed-data-Fourier}).
In this case, the planar equivalent-layer integral in the space domain (equation \ref{eq:eqlayer-integral}) assumes the form
\begin{equation}
	d^{o}(\mathbf{r}) \approx
	d(\mathbf{r})
	= \int\limits_{-\infty}^{\infty}\int\limits_{-\infty}^{\infty}
	p(\mathbf{r}'') \: \tfrac{\partial}{\partial z} \Psi(\mathbf{r}, \mathbf{r}'') \:\: d\mathcal{S}''
	\: , \quad z < z''_{0} \: .
	\label{eq:eqlayer-integral-Dirichlet-space}
\end{equation}
Note that the solution $p(\mathbf{r}'')$ obtained by computing the inverse Fourier transform of 
$\mathcal{F} [ p(\mathbf{r}'') ]$ (equation \ref{eq:p-Fourier-Dirichlet}) approximates the
downward continuation of the observed potential field to the layer plane.
From the physical point of view, Dirichlet's layers give rise to a predicted potential field $d(\mathbf{r})$ which is analogous
to the vertical component of the gravitational attraction due to a layer of monopoles
or the magnetic scalar potential due to a layer of vertically-magnetized dipoles.
Moreover, it can also be verified that in this case the planar equivalent-layer integral 
(equation \ref{eq:eqlayer-integral}) approximates the well-known \textit{upward continuation integral} 
\citep[e.g.,][]{skeels1947, peters1949, henderson-zietz1949, dean1958, roy1962, bhattacharyya1967, henderson1970},
which is traditionally derived from the solution of the Dirichlet’s problem on the
layer plane (equation \ref{eq:U-Dirichlet-Neumann} with $\ell = 0$). 
Dirichlet's equivalent layers are the basis of the pioneer method proposed by \cite{dampney1969} for processing gravity data.

The particular physical-property distribution $\tilde{p}(\mathbf{r}'')$ satisfying \ref{eq:eqlayer-integral-Dirichlet-space}
for a given observed potential field $d^{o}(\mathbf{r})$ leads to potential-field transformations (equations 
\ref{eq:transformed-field-continuous-layer} and \ref{eq:transformed-field-integral-Fourier}) defined by 
$\alpha_{t} = \nicefrac{k \, \alpha_{f}}{\alpha_{d}}$ (equation \ref{eq:alpha-t-generic}).
Dirichlet's equivalent layers are commonly used to perform upward/downward continuation of the observed potential 
field. In this case, $\alpha_{f} = \alpha_{d}$ and $\alpha_{t} = k$. 
Other potential-field transformations can also be accomplished from such equivalent layers, but they must be 
defined in terms of $\alpha_{f} \ne \alpha_{d}$ and $\alpha_{t} \ne k$.
In this case, however, defining the potential-field transformation in the space domain (equation \ref{eq:transformed-field-continuous-layer})
from its counterpart in the Fourier domain (equation \ref{eq:transformed-field-integral-Fourier}) is generally 
non-trivial.


\subsection{Equivalent layers having minimum phase changes}

\begin{definition}
	An equivalent layer has minimum phase changes (with respect to the spatial distribution of the true sources) if its
	physical-property distribution has a phase spectrum given by 
	\begin{equation}
		\Theta [ p(\mathbf{r}'') ] = \begin{cases}
			\Theta [ U(\mathbf{r}^{\dagger}) ] & , \quad 0 < k \le k_{\varepsilon} \\
			0 & , \quad k > k_{\varepsilon}
		\end{cases} \: ,
		\label{eq:p-phase-spectrum-U}
	\end{equation}
	where $z''_{0} > z^{\dagger}_{0}$ and $\Theta [ U(\mathbf{r}^{\dagger}) ]$ (Appendix B, equation \ref{eq:potential-U-phase-spectrum})
	is the phase spectrum of the potential $U(\mathbf{r})$ (equation \ref{eq:potential-U-space}).
	\label{def:eql-min-phase}
\end{definition}
\begin{corollary}
	\label{corollary:minimum-phase-change-layers}
	An equivalent layer having minimum phase changes (Definition \ref{def:eql-min-phase})
	depends only on the spatial distribution of the true sources and not on the potential field
	used to define the layer.
\end{corollary}
\begin{proof}
	It follows from equations \ref{eq:p-phase-spectrum-U} and \ref{eq:potential-U-phase-spectrum} that 
	$\mathcal{F} [ p(\mathbf{r}'') ]$ must be the product of a real function and
	$\mathcal{F} [ U(\mathbf{r}^{\dagger}) ]$, for $0 < k \le k_{\varepsilon}$, or $0$, for $k > k_{\varepsilon}$.
	From equations \ref{eq:potential-U-epsilon-Fourier-layer} and \ref{eq:p-Fourier}, we conclude that this condition is satisfied only for the case in which 
	$\alpha_{a} = \alpha_{r} \, \alpha_{d}$, where $\alpha_{r}$ is a real function satisfying 
	$\left| \alpha_{r} \right| \ne 0$ for $k > 0$.
	In this case, it can be verified 
	that $\mathcal{F} [ a(\mathbf{r}^{\dagger}, \mathbf{r}'') ]$ (equation \ref{eq:kernel-Fourier}) 
	assumes the form
	\begin{equation}
		\mathcal{F} [ a(\mathbf{r}^{\dagger}, \mathbf{r}'') ] =
		\alpha_{r} \, \alpha_{d} \, 
		\mathcal{F} [ \Psi(\mathbf{r}^{\dagger}, \mathbf{r}'') ] \: ,
		\quad \left| \alpha_{r} \right| \ne 0 \: ,
		\label{eq:kernel-Fourier-min-phase}
	\end{equation}
	where $\mathcal{F} [ \Psi(\mathbf{r}^{\dagger}, \mathbf{r}'') ]$ is given by equation \ref{eq:Psi-Fourier} 
	and $\mathcal{F} [ p(\mathbf{r}'') ]$ (equation \ref{eq:p-Fourier}) assumes the form
	\begin{equation}
		\mathcal{F} [ p(\mathbf{r}'') ] =
		\begin{cases}
			\tfrac{k}{\alpha_{r}} \, \mathcal{F} [ U(\mathbf{r}^{\dagger}) ] 
			e^{\, k \, (z''_{0} - z^{\dagger}_{0})} & , 
			\quad 0 < k \le k_{\varepsilon} \: , \quad \left| \alpha_{r} \right| \ne 0 \\
			0 & , \quad k > k_{\varepsilon}
		\end{cases} \: .
		\label{eq:p-Fourier-min-phase}
	\end{equation}
	Note that $\mathcal{F} [ p(\mathbf{r}'') ]$ (equation \ref{eq:p-Fourier-min-phase}) does not depend on $\alpha_{d}$.
	As a consequence, changes in the differential operator $\alpha_{d}$, which carries all the information related to
	the observed potential field, do not affect $p(\mathbf{r}'')$ (equation \ref{eq:eqlayer-integral}).
	This concludes the proof.
\end{proof}

The representation of equivalent layers having minimum phase changes (Definition \ref{def:eql-min-phase}) in the
space domain is usually non-trivial and depends on the particular form of the real function $\alpha_{r}$ 
(equations \ref{eq:kernel-Fourier-min-phase} and \ref{eq:p-Fourier-min-phase}).
Such layers give raise to potential-field transformations (equation \ref{eq:transformed-field-integral-Fourier})
defined by $\alpha_{t} = \alpha_{r} \, \alpha_{f}$ (equation \ref{eq:alpha-t-generic}), which may also
have non-trivial representations in the space domain.

%As a direct consequence of Corollary \ref{corollary:minimum-phase-change-layers}, we have that any 
%equivalent layer having minimum phase changes (Definition \ref{def:eql-min-phase}) 
%simultaneously approximates the scalar potential $U$ 
%and its partial derivatives at any horizontal plane $z^{\dagger}_{0} < z^{t}_{0}$ above the 
%sources according to the error bound \ref{eq:epsilon-bound-data}.
%Consequently, any potential-field transformation 
%(equations \ref{eq:transformed-field-continuous-layer} and \ref{eq:transformed-field-integral-Fourier})
%does not depend on the particular observed potential field used to define the equivalent layer as well.
%This can be shown by noting that, for a particular $\alpha_{a} = \tilde{\alpha}_{r} \, \tilde{\alpha}_{d}$,
%we have $\tilde{\alpha}_{t} = \left(\nicefrac{\tilde{\alpha}_{r}}{\tilde{\alpha}_{\Phi}}\right) \, \alpha_{f}$ 
%(equation \ref{eq:alpha-t-generic}), which does not
%depend on the $\tilde{\alpha}_{d}$ associated with the particular observed potential field.

\subsection{Neumann's equivalent layers}

\begin{definition}
	An equivalent layer is called Neumann's equivalent layer if the Fourier transform of its physical-property distribution has 
	the form
	\begin{equation}
		\mathcal{F} [ p(\mathbf{r}'') ] =
		\begin{cases}
			k \, \mathcal{F} [ U(\mathbf{r}^{\dagger}) ] 
			e^{\, k \, (z''_{0} - z^{\dagger}_{0})} & , 
			\quad 0 < k \le k_{\varepsilon} \: , \\
			0 & , \quad k > k_{\varepsilon} \: ,
		\end{cases} 
		\label{eq:p-Fourier-Neumann}
	\end{equation}
	where $z''_{0} > z^{\dagger}_{0}$.
	\label{def:eql-Neumann}
\end{definition}
\begin{lemma}
	All Neumann's equivalent layers (Definition \ref{def:eql-Neumann}) have 
	minimum phase changes according to Definition \ref{def:eql-min-phase}.
	\label{lemma:eql-Neumann}
\end{lemma}
\begin{proof}
	From equation \ref{eq:p-Fourier}, we have that the solution $\mathcal{F} [ p(\mathbf{r}'') ]$ 
	given by equation \ref{eq:p-Fourier-Neumann} to define a Neumann's equivalent layer
	(Definition \ref{def:eql-Neumann}) is satisfied 
	only for the case in which $\alpha_{a} = \alpha_{d}$. This condition implies that 
	$\mathcal{F} [ p(\mathbf{r}'') ]$ is a particular case of equation \ref{eq:p-Fourier-min-phase},
	with $\alpha_{r} = 1$, and consequently has minimum phase changes
	according to Definition \ref{def:eql-min-phase}.
\end{proof}

\subsubsection{General aspects}

Let a given observed potential field $d^{o}(\mathbf{r})$ be reproduced by a 
Neumann's equivalent layer (Definition \ref{def:eql-Neumann}). 
In this case, the planar equivalent-layer integral (equation \ref{eq:eqlayer-integral}) assumes the form
\begin{equation}
	d^{o}(\mathbf{r}) \approx
	d(\mathbf{r})
	= \int\limits_{-\infty}^{\infty}\int\limits_{-\infty}^{\infty}
	p(\mathbf{r}'') \: D_{d} \Psi(\mathbf{r}, \mathbf{r}'') \:\: d\mathcal{S}''
	\: , \quad z < z''_{0} \: ,
	\label{eq:eqlayer-integral-Neumann-space}
\end{equation}
where $D_{d}$ is the same space-domain differential operator associated with
the observed potential field data $d^{o}(\mathbf{r})$ (equation \ref{eq:observed-data-space}).
As a consequence, Neumann's equivalent layers (Definition \ref{def:eql-Neumann}) preserve the physical
relationship between the observed potential field and the true sources.
It means that, if the observed potential field $d^{o}(\mathbf{r})$ is the $x$-component of the 
gravitational attraction due to the true sources, the predicted potential field $d(\mathbf{r})$
represents the $x$-component of the gravitational attraction due to the planar equivalent layer, for example.
The same is valid for any other potential field.

Another interesting aspect of Neumann's layers is that, given a particular physical-property distribution 
$\tilde{p}(\mathbf{r}'')$ satisfying \ref{eq:eqlayer-integral-Neumann-space}, potential-field transformations 
are defined in the Fourier domain (equation \ref{eq:transformed-field-integral-Fourier}) by using 
$\alpha_{t} = \alpha_{f}$ (equation \ref{eq:alpha-t-generic}). Hence, potential-field transformations
(equation \ref{eq:transformed-field-continuous-layer}) assume the following form in the space domain:
\begin{equation}
	f(\mathbf{r})
	= \int\limits_{-\infty}^{\infty}\int\limits_{-\infty}^{\infty}
	\tilde{p}(\mathbf{r}'') \: D_{f} \Psi(\mathbf{r}, \mathbf{r}'') \:\: d\mathcal{S}''
	\: , \quad z < z''_{0} \: ,
	\label{eq:transformed-field-integral-Neumann-space}
\end{equation}
where $D_{f}$ is the space-domain representation of the differential operator $\alpha_{f}$ 
(equations \ref{eq:diff-operators-generic}, \ref{eq:kernel-transformation-Fourier} 
and \ref{eq:alpha-t-generic}) associated with the desired transformed field $f(\mathbf{r})$.

\subsubsection{Direct relationship with the true sources}

It follows immediately from equation \ref{eq:p-Fourier-Neumann} that, in the space domain, the physical-property 
distribution $p(\mathbf{r}'')$ of Neumann's equivalent layers (equation \ref{eq:eqlayer-integral-Neumann-space})
approximates the vertical derivative
$\tfrac{\partial}{\partial z} U(\mathbf{r}'')$ (equation \ref{eq:potential-U-1st-deriv-space}) on the 
layer plane $z''_{0}$.
For the cases in which $d^{o}(\mathbf{r})$ represents gravity data, the physical-property 
distribution $p(\mathbf{r}'')$ approximates the gravity disturbance (or vertical component of the gravitational
attraction) due to the true sources on the layer plane.
If $d^{o}(\mathbf{r})$ represents magnetic data produced by sources with constant total magnetization direction, 
then $p(\mathbf{r}'')$ approximates the pseudogravity 
anomaly \citep[][ p. 345]{blakely1996} due to the true sources on the layer plane.
In this case, $p(\mathbf{r}'')$ is strictly positive because it is defined in terms of a strictly positive 
function $\sigma(\mathbf{r}')$ (equation \ref{eq:potential-U-space}) representing the 
total-magnetization intensity at all points $\mathbf{r}' = (x', y', z')^{\top} \in \mathcal{V}$ 
within the sources.
This result is in perfect agreement with that deduced by \citet{reis_etal2020} for
the particular case in which the observed potential field $d^{o}(\mathbf{r})$ is
the approximated total-field anomaly (equation \ref{eq:potential-U-oblique-2nd-deriv-space}).

Since $p(\mathbf{r}'')$ (equation \ref{eq:eqlayer-integral-Neumann-space}) approximates 
$\tfrac{\partial}{\partial z} U(\mathbf{r}'')$ (equation \ref{eq:potential-U-1st-deriv-space}) on the 
layer plane $z''_{0}$, then it shares some common properties with the true sources.
The first is the total mass, for gravity sources, or total dipole moment, for magnetic sources with
constant total-magnetization direction.
For the cases in which the layer plane is located above the sources, $p(\mathbf{r}'')$ (equation \ref{eq:eqlayer-integral-Neumann-space})
reproduces $\tfrac{\partial}{\partial z} U(\mathbf{r}'')$ (equation \ref{eq:potential-U-1st-deriv-space}) exactly 
(equations \ref{eq:lim-kepsilon-p-Fourier}--\ref{eq:lim-kepsilon-predicted-data-space}) 
and its integration on the layer plane gives the total mass (or total dipole moment) according to the
Gauss' Theorem \citep[e.g.,][]{hammer1945, lafehr1965, bhattacharyya1967}. 
Furthermore, \citet{medeiros-silva1995a} show that the horizontal first-order moments, second-order moment $xy$, 
difference between the second-order moments $xx$ and $yy$, as well as the horizontal coordinates of the sources center of mass
(or center of dipole moment) of the true sources can also be estimated from the $\tfrac{\partial}{\partial z} U(\mathbf{r}'')$
that they produce on a horizontal plane above them.
These quantities can only be approximated for the case in which the layer crosses the sources because, in this case,
$p(\mathbf{r}'')$ does not reproduce $\tfrac{\partial}{\partial z} U(\mathbf{r}'')$ exactly.
The important aspect to be emphasized here is that the horizontal distribution of the physical-property distribution $p(\mathbf{r}'')$
(equation \ref{eq:eqlayer-integral-Neumann-space}) on Neumann's equivalent layers is equal to or approximates 
that associated with the true sources.

\subsubsection{Implications for gravity-field transformations}

In the context of gravity exploration, Lemma \ref{lemma:eql-Neumann} gives the 
theoretical foundation for the common approach to jointly process gravity-gradient data 
\cite[e.g.,][]{barnes_lumley2011} or to convert gravity disturbance data into the 
gravity-gradient tensor \cite[e.g.,][]{piauilino_etal2019_seg, oliveirajr_etal2023}, for example.
Such gravity-data transformations usually presume, without any proof, the existence of a single layer of 
monopoles that satisfactorily reproduces all components of the gravitational field.
This assumption, however, is a direct consequence of 
Corollary \ref{corollary:minimum-phase-change-layers} and Lemma \ref{lemma:eql-Neumann}.

For example, let us consider the common two-step approach to convert gravity-disturbance into gravity
gradient data. The first step consists in solving the planar equivalent-layer integral (equation \ref{eq:eqlayer-integral})
numerically for obtaining an estimated physical-property distribution $\tilde{p}(\mathbf{r}'')$ by considering that the
observed potential field is given by $d^{o}(\mathbf{r}) = \tfrac{\partial}{\partial z} U(\mathbf{r})$ (equation \ref{eq:potential-U-1st-deriv-space})
and that the kernel is $a(\mathbf{r}, \mathbf{r}'') = \tfrac{\partial}{\partial z} \Psi(\mathbf{r}, \mathbf{r}'')$ 
(equation \ref{eq:Psi}).
The second step consists in using the estimated $\tilde{p}(\mathbf{r}'')$ to compute the gradient components.
The $uv$-component of the gradient tensor is computed by setting $D_{t} \equiv \tfrac{\partial^{2}}{\partial u \, \partial v}$
(equations \ref{eq:potential-U-2nd-deriv-space} and \ref{eq:transformed-field-integral-Neumann-space}),
where $u, v \in \{ x, y, z \}$.

Note that the first step consists in defining a Neumann's equivalent layer (Definition \ref{def:eql-Neumann}) with 
$\alpha_{d} = \alpha_{a} \equiv \alpha(1, \hat{\mathbf{z}})$ 
(equations \ref{eq:observed-data-Fourier}, \ref{eq:diff-operators-generic} and \ref{eq:kernel-Fourier}) 
which leads to a particular physical-property distribution 
$\tilde{p}(\mathbf{r}'')$ having a Fourier transform given by \ref{eq:p-Fourier-Neumann}. The second step is equivalent to 
compute the potential-field transformations (equation \ref{eq:transformed-field-integral-Fourier}) defined in the Fourier domain 
by using $\alpha_{t} = \alpha_{f}$ (equation \ref{eq:alpha-t-generic}), with 
$\alpha_{f} \equiv \alpha(1, \hat{\mathbf{u}}) \, \alpha(1, \hat{\mathbf{v}})$ (equation \ref{eq:diff-operators-generic}), 
where the unit vectors $\hat{\mathbf{u}}$ and $\hat{\mathbf{v}}$ are given by the canonical base vectors (equation \ref{eq:canonical-base-vector}).
This Fourier-domain operator $\alpha_{f}$ is equivalent to the space-domain operator $D_{t} \equiv D(1, \hat{\mathbf{u}}) \, D(1, \hat{\mathbf{v}})$
(equation \ref{eq:diff-operators-generic}).

\subsubsection{Implications for magnetic-field transformations}

\begin{lemma}
	A planar equivalent layer of dipoles defined in terms of the magnetic scalar potential or a 
	component of the magnetic induction 
	field produced along a given constant direction by sources with a constant total-magnetization direction
	is a Neumann's equivalent layer (Definition \ref{def:eql-Neumann}) 
	if and only if its constant magnetization direction is equal to that defining the total magnetization of the sources.
	\label{lemma:eql-Neumann-dipoles}
\end{lemma}
\begin{proof}
	Again, we get from equations \ref{eq:p-Fourier} and \ref{eq:p-Fourier-Neumann} that all Neumann's 
	equivalent layers must be defined by $\alpha_{a} = \alpha_{d}$.
	We split this proof into two parts.
	
	In the first part, we consider that the observed potential field is the magnetic scalar potential. 
	In this case, the differential operator 
	$\alpha_{d} \equiv \alpha(1, \hat{\mathbf{v}})$ (equation \ref{eq:diff-operators-generic}),
	where $\hat{\mathbf{v}}$ is a unit vector with the same direction as that defining the
	sources total magnetization. Consider that 
	$\alpha_{a} \equiv \alpha(1, \hat{\mathbf{h}})$ (equation \ref{eq:diff-operators-generic}),
	where $\hat{\mathbf{h}}$ is a unit vector different from $\hat{\mathbf{v}}$. Suppose for sake of
	contradiction that $\alpha(1, \hat{\mathbf{v}}) = \alpha(1, \hat{\mathbf{h}})$. Then, we use 
	equation \ref{eq:diff-operators-generic} and the fact that the dot product is distributive to obtain 
	$(\hat{\mathbf{v}} - \hat{\mathbf{h}})^{\top} \mathbf{k} = 0$, which means that $\mathtt{i} k_{x}$, 
	$\mathtt{i} k_{y}$ and $k = \sqrt{k_{x}^{2} + k_{y}^{2}}$ are linearly dependent, a contradiction.
	This implies that $\hat{\mathbf{h}} = \hat{\mathbf{v}}$, completing the first part of our proof.
	
	In the second part of our proof, we consider that the observed potential field is a component
	of the magnetic induction field along a given constant direction defined by the unit vector $\hat{\mathbf{u}}$.
	Now, we have that 
	$\alpha_{d} \equiv \alpha(1, \hat{\mathbf{u}}) \, \alpha(1, \hat{\mathbf{v}})$ and
	$\alpha_{a} \equiv \alpha(1, \hat{\mathbf{u}}) \, \alpha(1, \hat{\mathbf{h}})$ 
	(equation \ref{eq:diff-operators-generic}). We consider that these two differential operators
	are equal to each other for sake of contradiction, which also leads to
	$(\hat{\mathbf{v}} - \hat{\mathbf{h}})^{\top} \mathbf{k} = 0$, another contradiction.
	This implies again that $\hat{\mathbf{h}} = \hat{\mathbf{v}}$, completing the second part of our proof.
\end{proof}

In the context of magnetic exploration methods, Lemmas \ref{lemma:eql-Neumann} and 
\ref{lemma:eql-Neumann-dipoles} contribute to a deeper understanding about the use 
of dipole layers to perform magnetic-field transformations.
Surprisingly, most of the previously published studies about this topic are limited to
the use of a planar layer of dipoles as a purely mathematical strategy for processing total-field 
anomaly data, as we can see in the comments of \citet{pedersen1989} on the work of \citet{bhattacharyya_chan1977},
for example.
Few studies \cite[][]{pedersen1991, li_etal2014, reis_etal2020} have drawn
attention to the problem of proving the existence of a planar layer of dipoles
that reproduces total-field anomaly data.
%By Lemmas \ref{lemma:eql-Neumann} and \ref{lemma:eql-Neumann-dipoles},
%we obtain that planar layers of dipoles having a constant 
%magnetization direction equal to that associated with the total-magnetization of the true sources
%are (i) Neumann's equivalent layers (Definition \ref{def:eql-Neumann}), (ii) intrinsically related 
%to the true magnetic sources and (iii) directly connected to the solution of the Neumann's problem on 
%a plane (equation \ref{eq:U-Neumann}). It means that such layers do not have a purely mathematical,
%but also a physical meaning.

By combining Lemmas \ref{lemma:eql-Neumann} and \ref{lemma:eql-Neumann-dipoles} with
Corollary \ref{corollary:minimum-phase-change-layers}, we obtain that
Neumann's equivalent layers of dipoles having a constant magnetization direction 
provides the theoretical basis for the common two-step strategy to compute the RTP in the space domain, 
from total-field anomaly data, according to the pioneer work of \citet{silva1986}.

The first step consists in solving the planar equivalent-layer integral (equation \ref{eq:eqlayer-integral})
numerically by considering that (i) the observed potential field 
is satisfactorily described by the directional derivative $d^{o}(\mathbf{r}) = \hat{\mathbf{u}}^{\top} \mathbf{H}_{U}(\mathbf{r}) \, \hat{\mathbf{v}}$ 
(equation \ref{eq:potential-U-oblique-2nd-deriv-space}), where $\hat{\mathbf{u}}$ defines the constant direction of the main geomagnetic
field at the study area and $\hat{\mathbf{v}}$ the constant total-magnetization direction of the sources, and (ii)
the kernel is $a(\mathbf{r}, \mathbf{r}'') = \hat{\mathbf{u}}^{\top} \mathbf{H}_{\Psi}(\mathbf{r}, \mathbf{r}'') \, \hat{\mathbf{v}}$,
where $\mathbf{H}_{\Psi}(\mathbf{r}, \mathbf{r}'')$ is the Hessian matrix \citep[e.g.,][ p. 158]{kaplan2002} of
function $\Psi(\mathbf{r}, \mathbf{r}'')$ (equation \ref{eq:Psi}).
The second step consists in using the estimated $\tilde{p}(\mathbf{r}'')$ to compute the RTP anomaly 
$\tfrac{\partial^{2}}{\partial z^{2}} U(\mathbf{r})$ (see the subsection \textit{Some relevant potential fields for geophysical applications}).
This is made by evaluating the potential-field transformation (equation \ref{eq:transformed-field-continuous-layer}) 
with $t(\mathbf{r}, \mathbf{r}'') = \tfrac{\partial^{2}}{\partial z^{2}} \Psi(\mathbf{r}, \mathbf{r}'')$ 
(equation \ref{eq:Psi}).

Note that the first step consists in defining a Neumann's equivalent layer (Definition \ref{def:eql-Neumann}) with 
$\alpha_{d} = \alpha_{a} \equiv \alpha(1, \hat{\mathbf{u}}) \, \alpha(1, \hat{\mathbf{v}})$ (equations \ref{eq:observed-data-Fourier}, 
\ref{eq:diff-operators-generic} and \ref{eq:kernel-Fourier}), which leads to a dipole layer having the same constant 
magnetization direction of the true sources and a particular physical-property distribution 
$\tilde{p}(\mathbf{r}'')$ with Fourier transform given by \ref{eq:p-Fourier-Neumann}. The second step is equivalent to  
compute the potential-field transformation (equation \ref{eq:transformed-field-integral-Fourier}) defined in the Fourier domain 
by using $\alpha_{t} = \alpha_{f}$ (equation \ref{eq:alpha-t-generic}), with 
$\alpha_{f} \equiv \alpha(2, \hat{\mathbf{z}})$ (equations \ref{eq:diff-operators-generic} and \ref{eq:canonical-base-vector}).
In this case, it can be verified from equations \ref{eq:diff-operators-generic} and \ref{eq:transformed-field-integral-Fourier3} that the 
transformed potential field $f(\mathbf{r}^{\ddagger})$ approximates the RTP anomaly $\tfrac{\partial^{2}}{\partial z^{2}} U(\mathbf{r}^{\ddagger})$ 
on the plane $z^{\ddagger}_{0}$.

It is important to underline two significant aspects that stands out from our theoretical approach to compute the RTP 
in terms of a Neumann's equivalent layer.
The first is that it has never been clearly defined in the pioneer work of \citet{silva1986}, nor in the succeeding studies 
\citep[e.g., ][]{leao-silva_1989, guspi_novara2009, oliveirajr_etal2013, li_etal2014, reis_etal2020, takahashi-etal2022} 
using the equivalent-layer technique. 
We do not presume, but mathematically prove the existence of a dipole layer that allows computing the RTP.
It is important to stress that \citet{silva1986} cites the work of \citet{roy1962} to justify that ``the magnetic 
field due to arbitrary sources can be exactly fitted by the magnetic field due to a horizontal equivalent
layer of dipoles''. 
Based on this, \citet{silva1986} proposed a method that (i) estimates the physical-property distribution on a dipole layer
having a constant magnetization direction equal to that of the true sources and (ii) sets the magnetization of this 
layer in the vertical direction to compute the RTP.
This two-step approach, that is widely accepted by the scientific community and we know to be correct from practice, is an empirical 
generalization of the \citeauthor{roy1962}'s (\citeyear{roy1962}) work.
\citeauthor{roy1962} did not discuss the existence of a dipole layer that fits the observed magnetic-field data,
but presented a general outline to deduce the solutions of the Dirichlet's and Neumann's problems on a plane 
(equation \ref{eq:U-Dirichlet-Neumann}).

The second significant aspect of our approach is that it also offers 
the theoretical link between the RTP in the space domain, via equivalent-layer technique, and 
in the Fourier domain, via classical approach \citep[e.g., ][ p. 331]{blakely1996}. 
This can be seen by first noting that the common two-step approach to compute the RTP via equivalent-layer technique is 
equivalent to a single-step approach defined by equation
\ref{eq:transformed-field-integral-Fourier-single-step2}, with $\alpha_{f} = \alpha(2, \hat{\mathbf{z}})$ 
(equations \ref{eq:diff-operators-generic}, \ref{eq:canonical-base-vector}, 
\ref{eq:transformed-field-integral-Fourier} and \ref{eq:alpha-t-generic}) and 
$\alpha_{d} \equiv \alpha(1, \hat{\mathbf{u}}) \, \alpha(1, \hat{\mathbf{v}})$ (equations \ref{eq:observed-data-Fourier} and 
\ref{eq:diff-operators-generic}), 
where $\hat{\mathbf{u}}$ defines the constant direction of the main geomagnetic
field at the study area and $\hat{\mathbf{v}}$ the constant total-magnetization direction of the sources.
Then, by considering that the plane $z^{\dagger}_{0}$ containing the observed data $d^{o}(\mathbf{r}^{\dagger})$ coincides with the 
plane $z^{\ddagger}_{0}$ where the transformed field $f(\mathbf{r}^{\ddagger})$ is computed, equation \ref{eq:transformed-field-integral-Fourier-single-step2} 
reduces to that commonly used in the literature \citep[e.g., ][ p. 331]{blakely1996} to perform the RTP in the 
Fourier domain.
This result gives the reason by which the equivalent-layer technique ultimately suffers the same instability as 
the Fourier-domain RTP operator at low latitudes, as have already been pointed out by \citet{li_etal2014}.


\subsection{Dipole layers with arbitrary constant magnetization direction}

Consider now that (i) the observed potential field 
is satisfactorily described by the directional derivative $d^{o}(\mathbf{r}) = \hat{\mathbf{u}}^{\top} \mathbf{H}_{U}(\mathbf{r}) \, \hat{\mathbf{v}}$ 
(equation \ref{eq:potential-U-oblique-2nd-deriv-space}), where $\hat{\mathbf{u}}$ defines the constant direction of the main geomagnetic
field at the study area and $\hat{\mathbf{v}}$ the constant total-magnetization direction of the sources, and (ii)
the kernel is $a(\mathbf{r}, \mathbf{r}'') = \hat{\mathbf{u}}^{\top} \mathbf{H}_{\Psi}(\mathbf{r}, \mathbf{r}'') \, \hat{\mathbf{h}}$,
where $\mathbf{H}_{\Psi}(\mathbf{r}, \mathbf{r}'')$ is the Hessian matrix \citep[e.g.,][ p. 158]{kaplan2002} of
function $\Psi(\mathbf{r}, \mathbf{r}'')$ (equation \ref{eq:Psi}) and $\hat{\mathbf{h}}$ is a constant unit vector not necessarily equal to 
$\hat{\mathbf{v}}$.
In this case, we get from equation \ref{eq:p-Fourier} that a solution of the planar equivalent layer 
can be obtained by using 
$\alpha_{d} \equiv \alpha(1, \hat{\mathbf{u}}) \, \alpha(1, \hat{\mathbf{v}})$ and 
$\alpha_{a} \equiv \alpha(1, \hat{\mathbf{u}}) \, \alpha(1, \hat{\mathbf{h}})$ 
(equation \ref{eq:diff-operators-generic}), so that the physical-property distribution assumes the following form 
in the Fourier domain:
\begin{equation}
	\mathcal{F} [ p(\mathbf{r}'') ] =
	\begin{cases}
		k \, \tfrac{\left(\mathbf{k}^{\top} \hat{\mathbf{v}} \right)}{\left(\mathbf{k}^{\top} \hat{\mathbf{h}} \right)} \,
		\mathcal{F} \left[ U(\mathbf{r}^{\dagger}) \right] 
		e^{\, k \, (z''_{0} - z^{\dagger}_{0})} & , 
		\quad 0 < k \le k_{\varepsilon} \\
		0 & , \quad k > k_{\varepsilon} 
	\end{cases} \: .
	\label{eq:p-Fourier-dipoles}
\end{equation}
This particular solution (equation \ref{eq:p-Fourier-dipoles}) was deduced by considering
that $\hat{\mathbf{u}}$ is a unit vector having the same constant direction of the main geomagnetic field at the study area.
However, it can be easily verified that this solution can also be deduced by considering that $\hat{\mathbf{u}}$ is one of the
canonical base vectors $\hat{\mathbf{x}}$, $\hat{\mathbf{y}}$ or $\hat{\mathbf{z}}$ (equation \ref{eq:canonical-base-vector}).
In this case, the observed potential field represents a Cartesian component of the magnetic induction field.
Additionally, we also deduce the same solution given by equation \ref{eq:p-Fourier-dipoles} if we consider that the 
observed potential field is the magnetic scalar potential. In this case, $\alpha_{d}$ and $\alpha_{a}$ are defined
without the term $\alpha(1, \hat{\mathbf{u}})$.
Hence, the planar equivalent layer formed by dipoles with physical-property distribution 
given by equation \ref{eq:p-Fourier-dipoles} simultaneously explain the magnetic scalar potential and
magnetic field components.

Another interesting aspect to be emphasized on the particular solution defined by equation \ref{eq:p-Fourier-dipoles} is that,
according to Lemma \ref{lemma:eql-Neumann-dipoles}, it 
does not necessarily lead to equivalent layers having minimum phase changes (Definition \ref{def:eql-min-phase}) owing to the complex ratio 
$\nicefrac{(\mathbf{k}^{\top} \hat{\mathbf{v}})}{(\mathbf{k}^{\top} \hat{\mathbf{h}})}$.
For the particular case in which $\hat{\mathbf{h}}$ coincides with the constant magnetization direction $\hat{\mathbf{v}}$ of the true sources,
then equation \ref{eq:p-Fourier-dipoles} reduces to that defining the physical-property distribution of a Neumann's layer 
(equation \ref{eq:p-Fourier-Neumann}), which has minimum phase changes according to Definition \ref{def:eql-min-phase}.

Note that the physical-property distribution $\mathcal{F} [ p(\mathbf{r}'') ]$ (equation \ref{eq:p-Fourier-dipoles}) associated 
with an arbitrarily magnetized dipole layer leads to potential-field transformations (equations 
\ref{eq:transformed-field-continuous-layer} and \ref{eq:transformed-field-integral-Fourier}) with 
$\alpha_{t}$ (equation \ref{eq:alpha-t-generic}) defined in terms of the complex ratio 
$\nicefrac{\alpha_{a}}{\alpha_{d}}$, where 
$\alpha_{d} \equiv \alpha(1, \hat{\mathbf{u}}) \, \alpha(1, \hat{\mathbf{v}})$ and 
$\alpha_{a} \equiv \alpha(1, \hat{\mathbf{u}}) \, \alpha(1, \hat{\mathbf{h}})$ 
(equation \ref{eq:diff-operators-generic}).
It means, for example, that the physical-property distribution $\mathcal{F} [ p(\mathbf{r}'') ]$ (equation \ref{eq:p-Fourier-dipoles}) 
cannot be used to compute the RTP by using the common approach with $\alpha_{t} = \alpha(2, \hat{\mathbf{z}})$ 
(equations \ref{eq:diff-operators-generic}, \ref{eq:canonical-base-vector} and \ref{eq:alpha-t-generic}).
This common approach can only be used if the constant magnetization direction $\hat{\mathbf{h}}$ coincides with 
that of the true sources $\hat{\mathbf{v}}$.
This fact is widely known by the geophysical community in practice and has long been assumed to be true
without any mathematical proof by several authors 
\citep[e.g., ][]{silva1986, leao-silva_1989, guspi_novara2009, oliveirajr_etal2013, takahashi-etal2022}. 
Our results offer the first theoretical explanation to this fact.

\subsection{Alternative equivalent layers}

\subsubsection{``Generalized equivalent sources'' proposed by \citet{cordell1992}}

Let the planar equivalent-layer integral (equation \ref{eq:eqlayer-integral}) be numerically solved by
using a kernel function $a(\mathbf{r}, \mathbf{r}'') = \Psi(\mathbf{r}, \mathbf{r}'')$ (equation \ref{eq:Psi}). 
Note that, in this case, the predicted potential field $d(\mathbf{r})$ 
represents the gravitational potential due to the equivalent layer.
The solution $p(\mathbf{r})$ obtained by this approach is defined from equation \ref{eq:p-Fourier}, in the Fourier domain, 
by setting $\alpha_{a} = 1$ (equation \ref{eq:diff-operators-generic}):
\begin{equation}
	\mathcal{F} [ p(\mathbf{r}'') ] =
	\begin{cases}
		k \, \alpha_{d} \, \mathcal{F} \left[ U(\mathbf{r}^{\dagger}) \right] 
		e^{\, k \, (z''_{0} - z^{\dagger}_{0})} & , 
		\quad 0 < k \le k_{\varepsilon} \\
		0 & , \quad k > k_{\varepsilon} 
	\end{cases} \: .
	\label{eq:p-Fourier-generalized-eqs}
\end{equation}
It is also important noting that 
this particular solution (equation \ref{eq:p-Fourier-generalized-eqs}) approximates the vertical derivative of the observed 
potential field on the layer plane and results in potential-field transformations (equations \ref{eq:transformed-field-continuous-layer} and 
\ref{eq:transformed-field-integral-Fourier}) defined by $\alpha_{t} = \nicefrac{\alpha_{f}}{\alpha_{d}}$ (equation \ref{eq:alpha-t-generic}).

This strategy, called \textit{generalized equivalent sources}, was proposed by \citet{cordell1992} and recently adapted by 
\citet{soler_uieda2021} to perform gravity-data processing. 
Those authors considered only the particular case in which observed potential field is gravity disturbance data 
$d^{o}(\mathbf{r}) = \tfrac{\partial}{\partial z} U(\mathbf{r})$ (equation \ref{eq:observed-data-space}), 
with $\alpha_{d} \equiv \alpha(1, \hat{\mathbf{z}})$ (equations \ref{eq:observed-data-Fourier}, 
\ref{eq:diff-operators-generic} and \ref{eq:canonical-base-vector}).
In this particular case, the solution given by equation \ref{eq:p-Fourier-generalized-eqs} 
(i) approximates $\tfrac{\partial^{2}}{\partial z^{2}} U(\mathbf{r})$ (equation \ref{eq:potential-U-2nd-deriv-space}) 
in the space domain, 
(ii) gives rise to equivalent layers having minimum phase changes (Definition \ref{def:eql-min-phase}), with 
$\alpha_{r} = \nicefrac{1}{k}$ (equation \ref{eq:p-Fourier-min-phase}), and 
(iii) results in potential-field transformations defined by $\alpha_{t} = \nicefrac{\alpha_{f}}{k}$ (equation \ref{eq:alpha-t-generic}).

\subsubsection{Alternative approach proposed by \citet{guspi_novara2009}}

\citet{guspi_novara2009} proposed an alternative two-step approach to compute the RTP via equivalent-layer technique. 
The first step consists in solving the planar 
equivalent-layer integral (equations \ref{eq:eqlayer-integral} and \ref{eq:eqlayer-integral-Fourier}) numerically 
by considering that (i) the observed potential field 
is satisfactorily described by the directional derivative 
$d^{o}(\mathbf{r}) = \hat{\mathbf{u}}^{\top} \mathbf{H}_{U}(\mathbf{r}) \, \hat{\mathbf{v}}$ 
(equation \ref{eq:potential-U-oblique-2nd-deriv-space}), where $\hat{\mathbf{u}}$ defines the constant direction of the main geomagnetic
field at the study area and $\hat{\mathbf{v}}$ the constant total-magnetization direction of the sources, and (ii)
the kernel is given by 
\begin{equation}
	a(\mathbf{r}, \mathbf{r}'') = \tfrac{\left( z''_{0} - z \right)}{2 \, \pi} \, 
	\ln{\Big[} \left( z''_{0} - z \right) + \| \mathbf{r}'' - \mathbf{r} \| {\Big]} - 
	\tfrac{\| \mathbf{r}'' - \mathbf{r} \|}{2 \, \pi} \: ,
	\label{eq:kernel-Fourier-Guspi}
\end{equation}
which represents the double integral of the function $\Psi(\mathbf{r}, \mathbf{r}'')$ (equation \ref{eq:Psi}) 
along the vertical direction $z$.
The second step of the alternative approach proposed by \citet{guspi_novara2009} computes the RTP anomaly 
by setting the potential-field transformation (equation \ref{eq:transformed-field-continuous-layer}) with 
$t(\mathbf{r}, \mathbf{r}'') = \Psi(\mathbf{r}, \mathbf{r}'')$ (equation \ref{eq:Psi}).

Note that the first step consists in obtaining an equivalent layer defined by 
$\alpha_{d} = (\mathbf{k}^{\top} \hat{\mathbf{u}}) (\mathbf{k}^{\top} \hat{\mathbf{v}})$ (equations \ref{eq:observed-data-Fourier} 
and \ref{eq:diff-operators-generic}) and 
$\alpha_{a} = \nicefrac{(\mathbf{k}^{\top} \hat{\mathbf{u}}) (\mathbf{k}^{\top} \hat{\mathbf{h}})}{k^{2}}$, 
where the term $\nicefrac{1}{k^{2}}$ represents the double integration along the vertical direction. In this case,  
the physical-property distribution $\mathcal{F} [ p(\mathbf{r}'') ]$ (equation \ref{eq:p-Fourier}) assumes the form:
\begin{equation}
	\mathcal{F} [ p(\mathbf{r}'') ] =
	\begin{cases}
		k^{3} \, \tfrac{\left(\mathbf{k}^{\top} \hat{\mathbf{v}} \right)}{\left(\mathbf{k}^{\top} \hat{\mathbf{h}} \right)} \,
		\mathcal{F} \left[ U(\mathbf{r}^{\dagger}) \right] 
		e^{\, k \, (z''_{0} - z^{\dagger}_{0})} & , 
		\quad 0 < k \le k_{\varepsilon} \\
		0 & , \quad k > k_{\varepsilon} 
	\end{cases} \: .
	\label{eq:p-Fourier-guspi}
\end{equation}
The second step is equivalent to compute the potential-field transformation (equation \ref{eq:transformed-field-integral-Fourier}) 
defined in the Fourier domain by using $\alpha_{t} = 1$ (equation \ref{eq:alpha-t-generic}).
In this case, it can be verified that the transformed potential field $\mathcal{F} [ f(\mathbf{r}^{\ddagger}) ]$ 
(equation \ref{eq:transformed-field-integral-Fourier}) actually approximates the RTP anomaly for the particular case 
in which $\mathcal{F} [ p(\mathbf{r}'') ]$ (equation \ref{eq:p-Fourier-guspi}) is defined by using $\hat{\mathbf{h}} = \hat{\mathbf{v}}$, 
i.e., for the case in which the constant magnetization direction of the equivalent layer coincides with that of the true sources.
This result shows the effectiveness of our approach in providing the theoretical background not only for classical 
equivalent-layer methods, but also for unconventional ones.
\section{Potential-field transformations via equivalent-layer technique}

Note that, given a particular equivalent layer $p_{\varepsilon}$ (Definition \ref{def:continuous-eqlayer}) 
defined as a square-integrable function by equation \ref{eq:p-Fourier}, in terms of specific $\alpha_{d}$ and $\alpha_{a}$, 
it is theoretically possible 
to compute any potential field transformation (equation \ref{eq:transformed-field-space} and \ref{eq:transformed-field-Fourier}).
This can be verified by first noting that applying the Fourier-convolution theorem to \ref{eq:transformed-field-continuous-layer} 
on a planar surface $\mathcal{S}^{\ddagger}$ at a constant vertical coordinate $z^{\ddagger}_{0} < z''_{0}$ gives
\begin{equation}
	\mathcal{F} [ f_{\varepsilon}(\mathbf{r}^{\ddagger}) ] =
	\mathcal{F} \left[ p_{\varepsilon}(\mathbf{r}'') \right]
	\mathcal{F} [ t(\mathbf{r}^{\ddagger}, \mathbf{r}'') ] \: , 
	\quad z^{\ddagger}_{0} < z''_{0} \: ,
	\label{eq:transformed-field-continuous-layer-Fourier}
\end{equation}
where
\begin{equation}
	\mathcal{F} [ t(\mathbf{r}^{\ddagger}, \mathbf{r}'') ] = \alpha_{t} \, \mathcal{F} [ \Psi(\mathbf{r}^{\ddagger}, \mathbf{r}'') ]
	\label{eq:kernel-t-Fourier}
\end{equation}
is defined in terms of a differential operator $\alpha_{t}$ acting on $\Psi$ (equation \ref{eq:Psi}) in the Fourier domain.
Then we set 
\begin{equation}
	\alpha_{t} = \left( \tfrac{\alpha_{a}}{\alpha_{d}} \right) \alpha_{f} \: ,
	\label{eq:alpha-t-generic}
\end{equation}
where $\alpha_{f}$ is the differential operator associated with the desired transformed field $f$ (equation \ref{eq:transformed-field-Fourier}).
In this case, we get from equations \ref{eq:Psi-Fourier}, \ref{eq:p-Fourier}  and \ref{eq:kernel-t-Fourier} 
that the transformed potential field $f_{\varepsilon}$ (equation \ref{eq:transformed-field-continuous-layer-Fourier}) 
is given by:
\begin{equation}
	\mathcal{F} [ f_{\varepsilon}(\mathbf{r}^{\ddagger}) ] =
	\alpha_{f} \, k \, \mathcal{F} \left[ U_{\varepsilon}(\mathbf{r}'') \right]
	\mathcal{F} [ \Psi(\mathbf{r}^{\ddagger}, \mathbf{r}'') ] \: , 
	\quad z^{\ddagger}_{0} < z''_{0} \: .
	\label{eq:transformed-field-continuous-layer-Fourier2}
\end{equation}
Finally, we use \ref{eq:Psi-Fourier}, \ref{eq:transformed-field-Fourier} and \ref{eq:potential-U-epsilon-Fourier-layer} to show that
\begin{equation}
	\mathcal{F} [ f_{\varepsilon}(\mathbf{r}^{\ddagger}) ] = 
	\begin{cases}
		\mathcal{F} [ f(\mathbf{r}^{\ddagger}) ] &
		, \quad 0 < k \le k_{\varepsilon} \\
		0 & , \quad k > k_{\varepsilon} 
	\end{cases} \: ,
	\label{eq:transformed-field-continuous-layer-1-Fourier}
\end{equation}
if the auxiliary potential $U_{\varepsilon}$ is defined by equation \ref{eq:potential-U-epsilon-Fourier-dagger-1}
or
\begin{equation}
	\mathcal{F} [ f_{\varepsilon}(\mathbf{r}^{\ddagger}) ] = 
	\tfrac{1}{1 + \mu_{\varepsilon} \, k^{2} \, e^{\, k \, (z''_{0} - z^{\dagger}_{0})}} \; 
	\mathcal{F} [ f(\mathbf{r}^{\ddagger}) ] \: ,
	\label{eq:transformed-field-continuous-layer-2-Fourier}
\end{equation}
if the auxiliary potential $U_{\varepsilon}$ is defined by equation \ref{eq:potential-U-epsilon-Fourier-dagger-2}.
In both cases, we can see that the transformed field $f_{\varepsilon}$ yielded by the equivalent layer 
is a smooth approximation of that produced by the true sources on the planar surface $\mathcal{S}^{\ddagger}$.

In the context of the equivalent-layer technique, a potential-field transformation is generally defined
in the space domain, by following two steps. The first consists in obtaining a particular layer $p_{\varepsilon}$ 
(Definition \ref{def:continuous-eqlayer}), which is subsequently used to compute a transformed field $f_{\varepsilon}$ via 
\ref{eq:transformed-field-continuous-layer}.
It can be verified, however, that this two-step approach is equivalent to the single-step approach based on the 
Fourier transform of the observed potential field (equation \ref{eq:field-transformation-Fourier}). 
We find by using \ref{eq:observed-data-Fourier} and \ref{eq:transformed-field-continuous-layer-Fourier2} that
\begin{equation}
	\mathcal{F} [ f_{\varepsilon}(\mathbf{r}^{\ddagger}) ] =
	\left( \tfrac{\alpha_{f}}{\alpha_{d}} \right) \, \mathcal{F} [ d_{\varepsilon}(\mathbf{r}^{\ddagger}) ]
	 \: .
	\label{eq:transformed-field-continuous-layer-single-step}
\end{equation}
Then, if we consider that the predicted field $d_{\varepsilon}$ (equation \ref{eq:predicted-data-continuous-layer-epsilon-Fourier})
if defined in terms of the auxiliary potential $U_{\varepsilon}$ given by equation \ref{eq:potential-U-epsilon-Fourier-dagger-1},
we get
\begin{equation}
	\mathcal{F} [ f_{\varepsilon}(\mathbf{r}^{\ddagger}) ] =
	\left( \tfrac{\alpha_{f}}{\alpha_{d}} \right) \, 
	\begin{cases}
		\mathcal{F} \left[ d^{o}(\mathbf{r}^{\dagger}) \right] 
		e^{\, k \, (z^{\ddagger}_{0} - z^{\dagger}_{0})} & , \quad 0 < k \le k_{\varepsilon} \\
		0 & , \quad k > k_{\varepsilon} 
	\end{cases} \: .
	\label{eq:transformed-field-continuous-layer-1-single-step}
\end{equation}
Alternatively, if we consider that the predicted field $d_{\varepsilon}$ (equation \ref{eq:predicted-data-continuous-layer-epsilon-Fourier})
if defined in terms of the auxiliary potential $U_{\varepsilon}$ given by equation \ref{eq:potential-U-epsilon-Fourier-dagger-2},
we get
\begin{equation}
	\mathcal{F} [ f_{\varepsilon}(\mathbf{r}^{\ddagger}) ] =
	\left( \tfrac{\alpha_{f}}{\alpha_{d}} \right) \, 
	\tfrac{1}{1 + \mu_{\varepsilon} \, k^{2} \, e^{\, k \, (z''_{0} - z^{\dagger}_{0})}} \;
	\mathcal{F} \left[ d^{o}(\mathbf{r}^{\dagger}) \right] \: .
	\label{eq:transformed-field-continuous-layer-2-single-step}
\end{equation}
Note that, for the particular case in which $z^{\ddagger}_{0}$ coincides with $z^{\dagger}_{0}$ and 
the truncating wavenumber $k_{\varepsilon}$ goes to infinity, equation 
\ref{eq:transformed-field-continuous-layer-1-single-step} reduces to \ref{eq:field-transformation-Fourier}.
Similarly, if $z^{\ddagger}_{0}$ coincides with $z^{\dagger}_{0}$ and $\mu_{\varepsilon}$ goes to zero,
equation \ref{eq:transformed-field-continuous-layer-2-single-step} also reduces to \ref{eq:field-transformation-Fourier}.
These results establish a theoretical link between potential-field transformations made via
equivalent-layer technique and those commonly made via Fourier transform of the observed data.
This link explains, for example, why potential-field transformations made via equivalent-layer technique may suffer from the same 
numerical problems generally associated with those computed via Fourier transform of the observed potential-field data.
\section{Equivalent layer classes}
\label{sec:eqlayer_classes}

\subsection{Dirichlet's equivalent layers}

\begin{definition}
	An equivalent layer $p_{\varepsilon}$ (Definition \ref{def:continuous-eqlayer}) is called Dirichlet's equivalent 
	layer if it has the following form in the Fourier domain
	\begin{equation}
		\mathcal{F} [ p_{\varepsilon}(\mathbf{r}'') ] =
		\alpha_{d} \, \mathcal{F} [ U_{\varepsilon}(\mathbf{r}'') ]  \: .
		\label{eq:p-Fourier-Dirichlet}
	\end{equation}
	\label{def:eql-Dirichlet}
\end{definition}
Note that Dirichlet's layers (Definition \ref{def:eql-Dirichlet}) are obtained from 
equation \ref{eq:p-Fourier} by setting $\alpha_{a} \equiv \alpha(1, \hat{\mathbf{z}}) = k$ 
(equation \ref{eq:diff-operators-generic}).
In this case, we can easily verify from equations \ref{eq:observed-data-Fourier},
\ref{eq:potential-U-epsilon-Fourier-layer-1} and \ref{eq:potential-U-epsilon-Fourier-layer-2} that they 
approximate the downward continuation of the observed potential field $d^{o}$ to the layer plane $\mathcal{S}''$.
It is important noting that such layers depend on $\alpha_{d}$, which 
carries information about the observed potential field (equation \ref{eq:observed-data-Fourier}).
In this case, the corresponding predicted field (equation \ref{eq:predicted-data-continuous-layer}) assumes the form
\begin{equation}
	d_{\varepsilon}(\mathbf{r})
	= \int\limits_{-\infty}^{\infty}\int\limits_{-\infty}^{\infty}
	p_{\varepsilon}(\mathbf{r}'') \: \tfrac{\partial}{\partial z} \Psi(\mathbf{r}, \mathbf{r}'') \:\: d\mathcal{S}''
	\: , \quad z < z''_{0} \: .
	\label{eq:predicted-data-Dirichlet-layer}
\end{equation}
From the physical point of view, this $d_{\varepsilon}$ is analogous
to the vertical component of the gravitational attraction due to a layer of monopoles
or the magnetic scalar potential due to a layer of vertically-magnetized dipoles.
Moreover, equation \ref{eq:predicted-data-Dirichlet-layer} approximates the well-known \textit{upward continuation integral} 
\citep[e.g.,][]{skeels1947, peters1949, henderson-zietz1949, dean1958, roy1962, bhattacharyya1967, henderson1970},
which is traditionally derived from the solution of the Dirichlet’s problem on the
layer plane (equation \ref{eq:U-Dirichlet-Neumann} with $\ell = 0$). 

Dirichlet's equivalent layers are the basis of the pioneer method proposed by \cite{dampney1969} for processing gravity data.
Such layers lead to potential-field transformations (equations 
\ref{eq:transformed-field-continuous-layer} and \ref{eq:transformed-field-continuous-layer-Fourier}) defined by 
$\alpha_{t} = \nicefrac{k \, \alpha_{f}}{\alpha_{d}}$ (equation \ref{eq:alpha-t-generic}).
Dirichlet's equivalent layers are commonly used to perform upward/downward continuation of the observed potential 
field, which is obtained by setting $\alpha_{f} = \alpha_{d}$, resulting in $\alpha_{t} = k$. 
Other potential-field transformations can also be accomplished from such equivalent layers, but they must be 
defined in terms of $\alpha_{f} \ne \alpha_{d}$.
In this case, however, defining the resulting transformed field $f_{\varepsilon}$ explicitly in the space domain 
(equation \ref{eq:transformed-field-continuous-layer}) 
from its counterpart in the Fourier domain (equation \ref{eq:transformed-field-continuous-layer-Fourier}) may be very difficult.


\subsection{Equivalent layers having minimum phase changes}

\begin{definition}
	An equivalent layer $p_{\varepsilon}$ (Definition \ref{def:continuous-eqlayer}) is said to have minimum 
	phase changes (with respect to the spatial distribution of the true sources) if its phase spectrum is given by
	\begin{equation}
		\Theta [ p_{\varepsilon}(\mathbf{r}'') ] = \Theta [ U_{\varepsilon}(\mathbf{r}^{\dagger}) ] \: ,
		\label{eq:p-phase-spectrum-U}
	\end{equation}
	where $\Theta [ U_{\varepsilon}(\mathbf{r}^{\dagger}) ]$ denotes the phase spectrum of $U_{\varepsilon}$ 
	(Appendix B, equations \ref{eq:potential-U-epsilon-1-phase-spectrum} and \ref{eq:potential-U-epsilon-2-phase-spectrum}).
	\label{def:eql-min-phase}
\end{definition}
\begin{proposition}
	\label{prop:minimum-phase-change-layers}
	An equivalent layer $p_{\varepsilon}$ having minimum phase changes (Definition \ref{def:eql-min-phase})
	depends only on the spatial distribution of the true sources and not on the particular observed 
	potential field $d^{o}$.
\end{proposition}
\begin{proof}
	It follows from equations \ref{eq:p-phase-spectrum-U} and definition of phase spectrum (Appendix B) that 
	$\mathcal{F} [ p_{\varepsilon}(\mathbf{r}'') ]$ must be the product of a real function and
	$\mathcal{F} [ U_{\varepsilon}(\mathbf{r}^{\dagger}) ]$.
	From equations \ref{eq:potential-U-epsilon-Fourier-layer} and \ref{eq:p-Fourier}, we can verify that 
	$\mathcal{F} [ p_{\varepsilon}(\mathbf{r}'') ]$ is in fact the product of 
	$\mathcal{F} [ U_{\varepsilon}(\mathbf{r}^{\dagger}) ]$ 
	and a term $\left( \nicefrac{\alpha_{d}}{\alpha_{a}} \right) k \, e^{\, k \, (z''_{0} - z^{\dagger}_{0})}$. 
	Since $\nicefrac{\alpha_{d}}{\alpha_{a}}$ is the only term that can be a complex function, we conclude that 
	the imposed condition can only be satisfied for the case in which $\alpha_{a} = \alpha_{r} \, \alpha_{d}$, 
	where $\alpha_{r}$ is a generic real function satisfying $\alpha_{r} \ne 0$ for $0 < k \le k_{\varepsilon}$.
	It can then be verified that, in this case, \ref{eq:p-Fourier} assumes the form
	\begin{equation}
		\mathcal{F} [ p_{\varepsilon}(\mathbf{r}'') ] =
		\tfrac{k}{\alpha_{r}} \, \mathcal{F} [ U_{\varepsilon}(\mathbf{r}^{\dagger}) ] 
		e^{\, k \, (z''_{0} - z^{\dagger}_{0})} \: ,
		\label{eq:p-Fourier-min-phase}
	\end{equation}
	which is the product of a real function and $\mathcal{F} [ U_{\varepsilon}(\mathbf{r}^{\dagger}) ]$.
	Since $\alpha_{r}$ does not depend on $\alpha_{d}$, $\mathcal{F} [ p_{\varepsilon}(\mathbf{r}'') ]$ does not 
	depend on $\alpha_{d}$ as well. This implies that the phase spectrum of $p_{\varepsilon}$ approximates that 
	of potential $U$ (Appendix B), thereby corresponding to an approximation 
	of the phase spectrum associated with the spatial distribution of the (shallow parts of the) true sources.
\end{proof}

Equivalent layers having minimum phase changes (Definition \ref{def:eql-min-phase}) may present intricate 
spatial domain representations, dictated by the particular form of the real function $\alpha_{r}$
(equation \ref{eq:p-Fourier-min-phase}).
Such layers give raise to potential-field transformations (equation \ref{eq:transformed-field-continuous-layer-Fourier})
defined by $\alpha_{t} = \alpha_{r} \, \alpha_{f}$ (equation \ref{eq:alpha-t-generic}), which may also
have non-trivial representations in the space domain.

\subsection{Neumann's equivalent layers}

\begin{definition}
	An equivalent layer $p_{\varepsilon}$ (Definition \ref{def:continuous-eqlayer}) is called Neumann's equivalent 
	layer if it has the following form in the Fourier domain
	\begin{equation}
		\mathcal{F} [ p_{\varepsilon}(\mathbf{r}'') ] =
		k \, \mathcal{F} [ U_{\varepsilon}(\mathbf{r}'') ] \: .
		\label{eq:p-Fourier-Neumann}
	\end{equation}
	\label{def:eql-Neumann}
\end{definition}
\begin{proposition}
	All Neumann's equivalent layers have minimum phase changes.
	\label{prop:eql-Neumann}
\end{proposition}
\begin{proof}
	Equation \ref{eq:p-Fourier-Neumann} is directly obtained from \ref{eq:p-Fourier}, by using 
	$\alpha_{a} = \alpha_{d}$. It can be easily verified that, in this case, 
	$\mathcal{F} [ p_{\varepsilon}(\mathbf{r}'') ]$ is a particular case of equation \ref{eq:p-Fourier-min-phase},
	with $\alpha_{r} = 1$. Consequently, Neumann's equivalent layers have minimum phase changes
	according to Definition \ref{def:eql-min-phase}.
\end{proof}

The predicted potential field $d_{\varepsilon}$ (equation \ref{eq:predicted-data-continuous-layer}) 
due to Neumann's equivalent layers (Definition \ref{def:eql-Neumann}) assumes the form
\begin{equation}
	d_{\varepsilon}(\mathbf{r})
	= \int\limits_{-\infty}^{\infty}\int\limits_{-\infty}^{\infty}
	p_{\varepsilon}(\mathbf{r}'') \: D_{d} \Psi(\mathbf{r}, \mathbf{r}'') \:\: d\mathcal{S}''
	\: , \quad z < z''_{0} \: ,
	\label{eq:predicted-data-Neumann-layer}
\end{equation}
where $D_{d}$ is the same space-domain differential operator associated with
the observed potential field data $d^{o}$ (equation \ref{eq:observed-data-space}).
As a consequence, Neumann's equivalent layers (Definition \ref{def:eql-Neumann}) preserve the physical
relationship between the observed potential field and the true sources.
It means that, if the observed potential field $d^{o}$ is the $x$-component of the 
gravitational attraction due to the true sources, the predicted potential field $d_{\varepsilon}$  
represents the $x$-component of the gravitational attraction due to the planar equivalent layer, for example.
The same is valid for any other potential field.

Any particular Neumann's layer $p_{\varepsilon}$ yields transformed potential 
fields (equation \ref{eq:transformed-field-continuous-layer-Fourier}) 
defined by $\alpha_{t} = \alpha_{f}$ (equation \ref{eq:alpha-t-generic}) in the Fourier domain. 
Hence, the corresponding transformed fields $f_{\varepsilon}$ in the space domain 
(equation \ref{eq:transformed-field-continuous-layer}) assume the simple form
\begin{equation}
	f_{\varepsilon}(\mathbf{r})
	= \int\limits_{-\infty}^{\infty}\int\limits_{-\infty}^{\infty}
	p_{\varepsilon}(\mathbf{r}'') \: D_{f} \Psi(\mathbf{r}, \mathbf{r}'') \:\: d\mathcal{S}''
	\: , \quad z < z''_{0} \: ,
	\label{eq:transformed-field-continuous-layer-Neumann}
\end{equation}
where $D_{f}$ is the space-domain representation of the differential operator $\alpha_{f}$ 
associated with the desired transformed field $f$ (equation \ref{eq:transformed-field-space}).

It follows immediately from equation \ref{eq:p-Fourier-Neumann} that Neumann's equivalent layers approximate 
the vertical derivative $\tfrac{\partial}{\partial z} U(\mathbf{r}'')$ (equation \ref{eq:potential-U-1st-deriv-space}) 
on the layer plane $z''_{0}$.
For the cases in which the observed potential field $d^{o}$ represents gravity data, 
$p_{\varepsilon}$ approximates the gravity disturbance (or vertical component of the gravitational
attraction) due to the true sources on the layer plane $\mathcal{S}''$.
If $d^{o}$ represents magnetic data produced by sources with constant total magnetization direction, 
then $p_{\varepsilon}$ approximates the pseudogravity 
anomaly \citep[][ p. 345]{blakely1996} due to the true sources on $\mathcal{S}''$.
In this case, $p_{\varepsilon}$ is strictly positive because it is defined in terms of a strictly positive 
function $\sigma(\mathbf{r}')$ (equation \ref{eq:potential-U-space}) representing the 
total-magnetization intensity at all points $\mathbf{r}' = (x', y', z')^{\top} \in \mathcal{V}$ 
within the sources.
This result is in perfect agreement with that deduced by \citet{reis_etal2020} for
the particular case in which the observed potential field $d^{o}$ is
the approximated total-field anomaly (equation \ref{eq:potential-U-oblique-2nd-deriv-space}).

It is widely known that integrating the potential field $\tfrac{\partial}{\partial z} U(\mathbf{r}'')$ 
(equation \ref{eq:potential-U-1st-deriv-space}) on a horizontal plane above the causative sources
gives their total mass (or total dipole moment), according to the
Gauss' Theorem \citep[e.g.,][]{hammer1945, lafehr1965, bhattacharyya1967}. 
Furthermore, \citet{medeiros-silva1995a} show that the horizontal first-order moments, second-order moment $xy$, 
difference between the second-order moments $xx$ and $yy$, as well as the horizontal coordinates of the sources center of mass
(or center of dipole moment) of the true sources can also be estimated from the $\tfrac{\partial}{\partial z} U(\mathbf{r}'')$
that they produce on a horizontal plane above them.
Thus, all these parameters of the true sources can be directly estimated from 
Neumann's equivalent layers $p_{\varepsilon}$ (Definition \ref{def:eql-Neumann}) because
they approximate the field $\tfrac{\partial}{\partial z} U(\mathbf{r}'')$ (equation \ref{eq:potential-U-1st-deriv-space}) 
on $\mathcal{S}''$.
%\section{Geophysical implications}

\subsection{Joint processing of gravity-gradient data}

The common approach to jointly process gravity-gradient data 
\cite[e.g.,][]{boggs-dransfield2004,barnes_lumley2011, martinez_li_2016} or to convert gravity disturbance data into the 
gravity-gradient tensor \cite[e.g.,][]{piauilino_etal2019_seg, oliveirajr_etal2023}, for example, 
usually presume, without any proof, the existence of a single layer of 
monopoles that satisfactorily reproduces the gravitational field and
preserves the physical relationship with the true sources.
This implies that the gravitational potential and its spatial derivatives attributed to the layer are presumed to 
accurately reproduce the corresponding field components produced by the true sources.
We provide the theoretical background for these potential-field transformations.

Let us consider, for example, the common two-step approach to convert gravity disturbance into gravity
gradient data (Table \partoneref{tbl:Table2}). The first step consists in obtaining a particular equivalent layer $p_{\varepsilon}$ 
by considering that the observed potential field $d^{o}$ is given by $\tfrac{\partial}{\partial z} U(\mathbf{r})$ 
(equation \partoneref{eq:potential-U-1st-deriv-space}) and that the kernel function 
$a(\mathbf{r}, \mathbf{r}'')$ is the vertical derivative of $\Psi$ (equation \partoneref{eq:Psi}).
The second step consists in using $p_{\varepsilon}$ to compute the gradient components.
The $uv$-component of the gradient tensor is computed by setting $D_{f} \equiv \tfrac{\partial^{2}}{\partial u \, \partial v}$
(equations \partoneref{eq:potential-U-2nd-deriv-space} and \partoneref{eq:transformed-field-continuous-layer-Neumann}),
where $u, v \in \{ x, y, z \}$.

Note that the first step consists in defining an equivalent layer 
by setting $\alpha_{d} = \alpha_{a} \equiv \alpha(1, \hat{\mathbf{z}})$ in equation \partoneref{eq:p-Fourier}.
According to Definition \partoneref{def:eql-Neumann}, this particular $p_{\varepsilon}$ is a Neumann's equivalent layer,
which naturally preserves the physical relationship with the true sources.
The second step of the common approach corresponds to 
compute transformed fields $f_{\varepsilon}$ (equation \partoneref{eq:transformed-field-continuous-layer-Fourier}) 
defined in the Fourier domain by using $\alpha_{t} = \alpha_{f}$ (equation \partoneref{eq:alpha-t-generic}), with 
$\alpha_{f} \equiv \alpha(1, \hat{\mathbf{u}}) \, \alpha(1, \hat{\mathbf{v}})$ (equation \partoneref{eq:diff-operators-generic}), 
where the unit vectors $\hat{\mathbf{u}}$ and $\hat{\mathbf{v}}$ are given by the canonical base vectors 
(equation \partoneref{eq:canonical-base-vector}).
We can easily get from \partoneref{eq:diff-operators-generic} that this Fourier-domain operator $\alpha_{f}$ is equivalent to the 
space-domain operator $D_{f} \equiv \tfrac{\partial^{2}}{\partial u \, \partial v}$.
Thus, the common two-step approach to convert gravity disturbance into gravity
gradient data via equivalent-layer technique can be theoretically explained as an application of 
a particular Neumann's equivalent layer (Definition \partoneref{def:eql-Neumann}).


\subsection{The existence of a planar dipole layer}


A predominant trend in previous studies utilizing a planar equivalent layer of dipoles for magnetic-field transformations 
is to treat the layer as a purely mathematical abstraction employed for processing total-field anomaly data (Table \partoneref{tbl:Table2}), 
a perspective exemplified by the remarks of \citet{pedersen1989} concerning the research of 
\citet{bhattacharyya_chan1977}.
Specifically, Pedersen pointed out that the first fundamental existence theorem for the Dirichlet's 
problem does not postulate that a general solution can be written as a layer of dipoles magnetized along its 
normal direction. In essence, Pedersen underscored that a dipole layer lacks physical significance.
In response, Chan affirmed that such a layer was employed in his work 
solely to address an integral equation, i.e., as a purely mathematical solution.

Since its inception in the 1960s, the equivalent-layer technique in applied geophysics has garnered 
relatively scant attention in terms of exploring the theoretical attributes of planar dipole layers 
capable of faithfully reproducing total-field anomaly data. 
This inquiry usually aims to discern the physical significance and possible link of such layers 
with boundary problems of potential theory.

Here, we leverage a theoretical approach to infer within the purview of applied geophysics, 
for the first time, the existence of planar dipole layers capable of accurately and simultaneously reproducing 
the magnetic scalar potential and induction field components along any direction, including that representing 
the approximated total-field anomaly (Table \partoneref{tbl:Table2}).

We first consider that (i) the observed potential field $d^{o}$ 
is satisfactorily described by the directional derivative $\hat{\mathbf{u}}^{\top} \mathbf{H}_{U}(\mathbf{r}) \, \hat{\mathbf{v}}$ 
(equation \partoneref{eq:potential-U-oblique-2nd-deriv-space}), where $\hat{\mathbf{u}}$ defines the constant direction of the main geomagnetic
field at the study area and $\hat{\mathbf{v}}$ the constant total-magnetization direction of the sources  (Table \partoneref{tbl:Table2}), and (ii)
the kernel is $a(\mathbf{r}, \mathbf{r}'') = \hat{\mathbf{u}}^{\top} \mathbf{H}_{\Psi}(\mathbf{r}, \mathbf{r}'') \, \hat{\mathbf{h}}$,
where $\mathbf{H}_{\Psi}$ is the Hessian matrix \citep[e.g.,][ p. 158]{kaplan2002} of
function $\Psi$ (equation \partoneref{eq:Psi}) and $\hat{\mathbf{h}}$ is a constant unit vector not necessarily equal to 
$\hat{\mathbf{v}}$.
In this case, we use $\alpha_{d} \equiv \alpha(1, \hat{\mathbf{u}}) \, \alpha(1, \hat{\mathbf{v}})$ and 
$\alpha_{a} \equiv \alpha(1, \hat{\mathbf{u}}) \, \alpha(1, \hat{\mathbf{h}})$ 
(equation \partoneref{eq:diff-operators-generic}) in equation \partoneref{eq:p-Fourier} to obtain the following equivalent layer 
in the Fourier domain:
\begin{equation}
	\mathcal{F} [ p_{\varepsilon}(\mathbf{r}'') ] =
	k \, \tfrac{\left(\mathbf{k}^{\top} \hat{\mathbf{v}} \right)}{\left(\mathbf{k}^{\top} \hat{\mathbf{h}} \right)} \,
	\mathcal{F} \left[ U_{\varepsilon}(\mathbf{r}'') \right] \: .
	\label{eq:p-Fourier-dipoles}
\end{equation}
This particular equivalent layer represents a planar layer of dipoles with constant magnetization direction
$\hat{\mathbf{h}}$ and was deduced by considering
that $\hat{\mathbf{u}}$ is a unit vector having the same constant direction of the main geomagnetic field at the study area.
However, it can be easily verified that the same equivalent layer can also be deduced by considering that $\hat{\mathbf{u}}$ is one of the
canonical base vectors $\hat{\mathbf{x}}$, $\hat{\mathbf{y}}$ or $\hat{\mathbf{z}}$ (equation \partoneref{eq:canonical-base-vector}).
In this case, the observed potential field represents a Cartesian component of the magnetic induction field  (Table \partoneref{tbl:Table2}).
It can also be easily verified that the same $p_{\varepsilon}$ given by equation 
\ref{eq:p-Fourier-dipoles} can be equally deduced if we consider that the 
observed potential field is the magnetic scalar potential  (Table \partoneref{tbl:Table2}). 
In this case, $\alpha_{d}$ and $\alpha_{a}$ are defined
without the term $\alpha(1, \hat{\mathbf{u}})$. 

The main conclusion obtained from the reasoning presented above is that a planar equivalent layer of dipoles defined according 
to equation \ref{eq:p-Fourier-dipoles} simultaneously 
reproduces the magnetic scalar potential and any magnetic induction field component (including the approximated total-field anomaly) 
due to sources with constant total-magnetization direction defined by $\hat{\mathbf{v}}$, even if the layer has a constant magnetization 
direction defined by a unit vector $\hat{\mathbf{h}} \ne \hat{\mathbf{v}}$. 
This theoretical result explains, for example, why we can use a planar equivalent layer of dipoles with arbitrary constant 
magnetization direction $\hat{\mathbf{h}}$ to compute the magnetic induction field components or the amplitude of the magnetic induction field due to the true sources 
without presuming any information about the constant direction of their total magnetization 
\citep[e.g,][]{li_etal_2021}.


\subsection{Computing the RTP by using a dipole layer}


According to the pioneer work of \citet{silva1986}, which is widely followed in applied geophysics, 
the RTP in the space domain is calculated from total-field anomaly data (Table \partoneref{tbl:Table2}) in two steps.
The first consists in estimating the magnetic-moment intensity distribution on a planar layer 
of dipoles (or doublets) with constant magnetization direction equal to that of the true sources.
The second step consists in computing the $z$-component of the magnetic induction field due 
to the estimated layer by setting its magnetization along the vertical direction.
It is usually verified a posteriori that this computed field approximates the true RTP anomaly (Table \partoneref{tbl:Table2}).

We call attention to the fact that \citet{silva1986} cited the work of \citet{roy1962} to justify that ``the magnetic 
field due to arbitrary sources can be exactly fitted by the magnetic field due to a horizontal equivalent
layer of dipoles''. 
However, \citeauthor{roy1962} neither discussed the existence of a dipole layer that fits the observed total-field 
anomaly data nor its use to retrieve the RTP anomaly, but 
presented a general outline to deduce the solutions of the Dirichlet's and Neumann's problems on a plane 
located above the sources (equation \partoneref{eq:U-Dirichlet-Neumann}).
In other words, the two-step approach proposed by \citet{silva1986}, which is widely accepted by the scientific community 
and we know to be correct from practice 
\citep[e.g., ][]{leao-silva_1989, guspi_novara2009, oliveirajr_etal2013, li_etal2014, reis_etal2020, takahashi-etal2022, piauilino_etal2025}, 
is an empirical generalization of the \citeauthor{roy1962}'s (\citeyear{roy1962}) work. 
It is also important to point out that the mathematical equivalence between the space-domain approach to compute the 
RTP via equivalent-layer technique and the classical approach directly using the Fourier transform of the observed 
total-field anomaly data \citep[e.g., ][ p. 331]{blakely1996} have never been clearly established in the literature.
Here, we provide a general approach aiming at filling these theoretical gaps.

From a theoretical point of view, the two-step approach to compute the RTP in the space domain 
consists in first setting a planar layer of dipoles according to equation \ref{eq:p-Fourier-dipoles}, 
by considering the unit vector $\hat{\mathbf{h}}$ equal to $\hat{\mathbf{v}}$, which defines 
the constant total-magnetization direction of the true sources.
In this case, the ratio $\nicefrac{\left(\mathbf{k}^{\top} \hat{\mathbf{v}} \right)}{\left(\mathbf{k}^{\top} \hat{\mathbf{h}} \right)} = 1$ 
and we can easily verify that the $\mathcal{F} [ p_{\varepsilon}(\mathbf{r}'') ]$ given by 
equation \ref{eq:p-Fourier-dipoles} assumes the form of a Neumann's equivalent layer (Definition \partoneref{def:eql-Neumann}).
The second step to compute the RTP consists in using the predefined Neumann's layer to compute 
the transformed potential field $f_{\varepsilon}$ (equation \partoneref{eq:transformed-field-continuous-layer}) defined in the Fourier domain 
by using $\alpha_{t} = \alpha_{f}$ (equation \partoneref{eq:alpha-t-generic}), with 
$\alpha_{f} \equiv \alpha(2, \hat{\mathbf{z}})$ (equations \partoneref{eq:diff-operators-generic} and \partoneref{eq:canonical-base-vector}).
In this case, it can be verified from equations \partoneref{eq:diff-operators-generic} and \partoneref{eq:transformed-field-continuous-layer-single-step} 
that the computed transformed field $f_{\varepsilon}$ (equation \partoneref{eq:transformed-field-continuous-layer}) 
approximates the RTP anomaly  (Table \partoneref{tbl:Table2}). 
As we have just shown here, for the first time, the classical approach of \citet{silva1986} to compute the RTP in the 
space domain can be formally explained in terms of a Neumann's equivalent layer (Definition \partoneref{def:eql-Neumann}). 

We can also investigate the theoretical link between the RTP in the space domain, 
via equivalent-layer technique, and in the Fourier domain, via classical approach \citep[e.g., ][ p. 331]{blakely1996}. 
As we have already discussed, the transformed field $f_{\varepsilon}$ computed via \partoneref{eq:transformed-field-continuous-layer-Fourier} 
from a previously defined equivalent layer $p_{\varepsilon}$ can also be directly 
computed from the observed potential field $d^{o}$ (equation \partoneref{eq:transformed-field-continuous-layer-1-single-step}
or \partoneref{eq:transformed-field-continuous-layer-2-single-step}), 
in terms of the ratio $\nicefrac{\alpha_{f}}{\alpha_{d}}$.
To compute the RTP, we set $\alpha_{d} \equiv \alpha(1, \hat{\mathbf{u}}) \, \alpha(1, \hat{\mathbf{v}})$ 
and $\alpha_{f} \equiv \alpha(2, \hat{\mathbf{z}})$ (equation \partoneref{eq:diff-operators-generic}).
By using these specific $\alpha_{d}$ and $\alpha_{f}$, setting $z^{\ddagger}_{0} = z^{\dagger}_{0}$ and 
(i) letting $k_{\varepsilon}$ goes to infinity in equation \partoneref{eq:transformed-field-continuous-layer-1-single-step} or
(ii) letting $\mu_{\varepsilon}$ goes to zero in equation \partoneref{eq:transformed-field-continuous-layer-2-single-step}, we get 
\begin{equation}
	\mathcal{F} [ f_{\varepsilon}(\mathbf{r}^{\dagger}) ] =
	\tfrac{k^{2}}{\left(\mathbf{k}^{\top} \hat{\mathbf{u}} \right) \, \left(\mathbf{k}^{\top} \hat{\mathbf{v}} \right)} \, 
	\mathcal{F} \left[ d^{o}(\mathbf{r}^{\dagger}) \right] \: ,
	\label{eq:transformed-field-continuous-layer-single-step-RTP}
\end{equation}
which represents the
classical approach to compute the RTP in the Fourier domain \citep[e.g., ][ p. 331]{blakely1996}.
We stress two important aspects of this result. The first concerns the theoretical link 
between the space- and Fourier-domain approaches to compute the RTP. To the authors knowledge, this link has never been fully 
detailed in the literature.
The second important aspect of our theoretical approach is that it provides the 
reason by which the equivalent-layer technique ultimately suffers the same instability as 
the Fourier-domain RTP operator at low latitudes. In this case, the estimated equivalent layer of dipoles 
presents striations, as have already been pointed out by \citet{li_etal2014}. 
In the Fourier domain, these artifacts are directly related to the term 
$\nicefrac{\left(\mathbf{k}^{\top} \hat{\mathbf{v}} \right)}{\left(\mathbf{k}^{\top} \hat{\mathbf{h}} \right)}$ 
in equation \ref{eq:p-Fourier-dipoles}. This term produces, in the equivalent layer of dipoles, a 
numerical instability similar to that caused by the term 
$\nicefrac{k^{2}}{\left(\mathbf{k}^{\top} \hat{\mathbf{u}} \right) \, \left(\mathbf{k}^{\top} \hat{\mathbf{v}} \right)}$ in equation 
\ref{eq:transformed-field-continuous-layer-single-step-RTP},
which is the base of the classical approach to compute the RTP in the Fourier domain \citep[e.g., ][ p. 331]{blakely1996}.


\subsection{Generalized equivalent sources}

Let us define an equivalent layer $p_{\varepsilon}$ by considering the kernel function 
$a(\mathbf{r}, \mathbf{r}'') = \Psi(\mathbf{r}, \mathbf{r}'')$ (equation \partoneref{eq:Psi}). 
Note that, in this case, 
the corresponding predicted potential field $d_{\varepsilon}$ (equation \partoneref{eq:predicted-data-continuous-layer}) 
is analogous to the gravitational potential due to the equivalent layer.
Besides, according to equations \partoneref{eq:p-Fourier}--\partoneref{eq:kernel-a-Fourier}, we have that 
$\alpha_{a} = 1$ and 
\begin{equation}
	\mathcal{F} [ p_{\varepsilon}(\mathbf{r}'') ] =
	k \, \alpha_{d} \, \mathcal{F} \left[ U_{\varepsilon}(\mathbf{r}'') \right]  \: ,
	\label{eq:p-Fourier-generalized-eqs}
\end{equation}
which approximates the vertical derivative of the observed potential field on the layer plane. 
Such layers produce potential-field transformations (equations \partoneref{eq:transformed-field-continuous-layer} and 
\partoneref{eq:transformed-field-continuous-layer-Fourier}) defined by $\alpha_{t} = \nicefrac{\alpha_{f}}{\alpha_{d}}$ (equation \partoneref{eq:alpha-t-generic}).

This strategy, called \textit{generalized equivalent sources}, was proposed by \citet{cordell1992} and recently adapted by 
\citet{soler_uieda2021} to perform gravity-data processing. 
Those authors considered only the particular case in which the observed potential field $d^{o}$ is gravity disturbance data 
$\tfrac{\partial}{\partial z} U(\mathbf{r})$  (Table \partoneref{tbl:Table2}), 
with $\alpha_{d} \equiv \alpha(1, \hat{\mathbf{z}}) = k$ (equations \partoneref{eq:observed-data-Fourier}, 
\partoneref{eq:diff-operators-generic} and \partoneref{eq:canonical-base-vector}).
In this particular case, equation \ref{eq:p-Fourier-generalized-eqs} leads to an equivalent layer that 
(i) approximates $\tfrac{\partial^{2}}{\partial z^{2}} U(\mathbf{r})$ (equation \partoneref{eq:potential-U-2nd-deriv-space}) 
in the space domain, 
(ii) has minimum phase changes (Definition \partoneref{def:eql-min-phase}) and 
(iii) results in potential-field transformations defined by $\alpha_{t} = \nicefrac{\alpha_{f}}{k}$ (equation \partoneref{eq:alpha-t-generic}).

\subsection{Alternative approach to compute the RTP}

\citet{guspi_novara2009} proposed an alternative two-step approach to compute the RTP via equivalent-layer technique. 
The first step consists in obtaining an equivalent layer by considering that (i) the observed potential field $d^{o}$
is satisfactorily described by the directional derivative 
$\hat{\mathbf{u}}^{\top} \mathbf{H}_{U}(\mathbf{r}) \, \hat{\mathbf{v}}$ 
(equation \partoneref{eq:potential-U-oblique-2nd-deriv-space}), where $\hat{\mathbf{u}}$ defines the constant direction of the main geomagnetic
field at the study area and $\hat{\mathbf{v}}$ the constant total-magnetization direction of the sources  (Table \partoneref{tbl:Table2}), and (ii)
the kernel function is given by 
\begin{equation}
	a(\mathbf{r}, \mathbf{r}'') = \tfrac{\left( z''_{0} - z \right)}{2 \, \pi} \, 
	\ln{\Big[} \left( z''_{0} - z \right) + \| \mathbf{r}'' - \mathbf{r} \| {\Big]} - 
	\tfrac{\| \mathbf{r}'' - \mathbf{r} \|}{2 \, \pi} \: ,
	\label{eq:kernel-Fourier-Guspi}
\end{equation}
which represents the double integral of the function $\Psi$ (equation \partoneref{eq:Psi}) 
along the vertical direction $z$.
The second step of the alternative approach proposed by \citet{guspi_novara2009} computes the RTP anomaly 
(Table \partoneref{tbl:Table2}) by setting $f_{\varepsilon}$ (equation \partoneref{eq:transformed-field-continuous-layer}) with 
$t(\mathbf{r}, \mathbf{r}'') = \Psi(\mathbf{r}, \mathbf{r}'')$ (equation \partoneref{eq:Psi}).

Note that the first step consists in obtaining an equivalent layer $p_{\varepsilon}$ defined by 
$\alpha_{d} = (\mathbf{k}^{\top} \hat{\mathbf{u}}) (\mathbf{k}^{\top} \hat{\mathbf{v}})$ (equations \partoneref{eq:observed-data-Fourier} 
and \partoneref{eq:diff-operators-generic}) and 
$\alpha_{a} = \nicefrac{(\mathbf{k}^{\top} \hat{\mathbf{u}}) (\mathbf{k}^{\top} \hat{\mathbf{h}})}{k^{2}}$, 
where the term $\nicefrac{1}{k^{2}}$ represents the double integration along the vertical direction. In this case,  
$p_{\varepsilon}$ assumes the following form in the Fourier domain:
\begin{equation}
	\mathcal{F} [ p_{\varepsilon}(\mathbf{r}'') ] =
	k^{3} \, \tfrac{\left(\mathbf{k}^{\top} \hat{\mathbf{v}} \right)}{\left(\mathbf{k}^{\top} \hat{\mathbf{h}} \right)} \,
	\mathcal{F} \left[ U_{\varepsilon}(\mathbf{r}'') \right] \: .
	\label{eq:p-Fourier-guspi}
\end{equation}
The second step consists in computing $f_{\varepsilon}$ in the Fourier domain (equation \partoneref{eq:transformed-field-continuous-layer-Fourier}) 
by using $\alpha_{t} = 1$ (equation \partoneref{eq:alpha-t-generic}).
In this case, it can be verified that the transformed potential field $\mathcal{F} [ f_{\varepsilon}(\mathbf{r}^{\ddagger}) ]$ 
(equation \partoneref{eq:transformed-field-continuous-layer-Fourier}) actually approximates the RTP anomaly  (Table \partoneref{tbl:Table2}) 
for the particular case in which $\hat{\mathbf{h}} = \hat{\mathbf{v}}$, 
i.e., for the case in which the constant magnetization direction of the equivalent layer coincides with that of the true sources.
This result shows the effectiveness of our approach in providing the theoretical background not only for classical 
equivalent-layer methods, but also for unconventional ones.
\section{Conclusions}
\label{sec:conclusions}

We furnish the theoretical background for potential-field transformations
that have been empirically made in space domain by using equivalent layers since 
this technique was introduced in applied geophysics during the 1960s.
Our goal here is to replace the widely accepted equivalent-layer principle, which 
refers to the empirical generalizations
by which most of the studies found in the geophysical literature justify the existence of
equivalent layers allowing the computation of potential-field transformations.
Specifically, we have presented the planar equivalent-layer theorem to state
the existence of planar and continuous equivalent layers that 
reproduce the potential-field data exactly when they reside strictly in the source-free 
region above the sources, or within a predefined error level when they intersect the true sources.
We have proven this mathematical statement by modifying the solutions of the 
Dirichlet's and Neumann's problems on planar surfaces that may cross the sources.
We also establish how a particular equivalent layer can be used to perform a particular 
potential-field transformation.
Furthermore, we define the theoretical link between potential-field transformations made 
via the equivalent-layer technique and those directly computed from the Fourier transform of the observed 
potential-field data. 
In essence, we have filled the theoretical gaps in most of the existing equivalent-layer 
methods and have also shown that a planar equivalent layer is not a purely mathematical abstraction,
but carries physical meaning and is directly related to boundary value problems of potential 
theory.
We contribute to the general understanding on why the technique has long been 
successfully applied to perform potential-field transformations. 
Our work opens the door for future theoretical investigations about potential-field 
transformations via the equivalent-layer technique, as well as for the relationship between 
planar equivalent layers and the true sources.

\section{Acknowledgments}

This study was financed by the Brazilian agencies FAPERJ (grant nos. E-26 202.729/2018 and 26/202.582/2019) 
and CNPq (grant nos. 315768/2020-7 and 309624/2021-5). The authors thank the editors and reviewers 
who contributed to this work.

\clearpage

%\include{algorithms}
\tabl{Table1}{Table of symbols.
}{
	\begin{center}
		\renewcommand{\arraystretch}{0.4}
		\begin{tabular}{
				>{\smaller\raggedleft\arraybackslash}m{0.18\textwidth}
				>{\tiny\setlength{\baselineskip}{0.5\baselineskip}\raggedright\arraybackslash}m{0.32\textwidth}
				>{\smaller\raggedleft\arraybackslash}m{0.18\textwidth}
				>{\tiny\setlength{\baselineskip}{0.5\baselineskip}\raggedright\arraybackslash}m{0.32\textwidth}
			}
			{\small \textbf{Symbol}} & 
			{\small \textbf{Definition}} & 
			{\small \textbf{Symbol}} & 
			{\small \textbf{Definition}} \\
			\hline
			$\mathcal{V}'_{\nu}$ & 
			Regular region: Open, connected subset of $\mathbb{R}^{3}$, with boundary defined by a closed and compact 
			surface $\partial \mathcal{V}'_{\nu}$ free of double points \citep[][ p. 22]{freeden_gerhards2012}. &
			$\overline{\mathcal{V}'}$ & 
			Closed region defined as the union of a finite collection of $N_{\nu}$ disjoint regular regions $\mathcal{V}'_{\nu}$ 
			and their respective boundaries $\partial\mathcal{V}'_{\nu}$, i.e., 
			$\overline{\mathcal{V}'} = \bigcup_{\nu = 1}^{N_{\nu}} \overline{\mathcal{V}'_{\nu}}$, where
			$\overline{\mathcal{V}'_{\nu}} = \mathcal{V}'_{\nu} \cup \partial\mathcal{V}'_{\nu}$.
			We consider that $\overline{\mathcal{V}'}$ contains the sources. \\
			$z^{t}_{0}$ and $\mathbf{r}'$ &
			Constant vertical coordinate defining the top of the sources and 
			generic point $\mathbf{r}' = (x', y', z')^{\top}$ within the sources, i.e, $\mathbf{r}' \in \overline{\mathcal{V}'}$. &
			$\sigma(\mathbf{r}')$ & 
			Piecewise continuous, integrable and bounded function that satisfies some Hölder condition \citep[][ p. 152]{kellogg1967} 
			at all points $\mathbf{r}' \in \overline{\mathcal{V}'}$ and is zero at all points $\mathbf{r} \notin \overline{\mathcal{V}'}$.
			This function represents the density (contrast) or total-magnetization intensity within the sources. \\
			$\mathcal{S}''$ and $\mathbf{r}''$ & 
			The first variable represents a horizontal surface located at a constant vertical coordinate $z''_{0}$ that may cross 
			the sources. This surface is used to define all equivalent layers in the present work. 
			The second variable defines a generic point $\mathbf{r}'' = (x'', y'', z''_{0})^{\top} \in \mathcal{S}''$. &
			$\mathbf{r}$ &
			Generic point $\mathbf{r} = (x, y, z)^{\top}$ outside the sources ($\mathbf{r} \notin \overline{\mathcal{V}'}$). 
			It is also used to define generic observation points with potential-field data located above the sources and the surface 
			$\mathcal{S}''$, i.e., with vertical coordinates satisfying $z < z^{t}_{0}$ and $z < z''_{0}$. \\
			$\mathcal{S}^{\dagger}$ and $\mathcal{S}^{\ddagger}$ & 
			Horizontal surfaces located at constant vertical coordinates $z^{\dagger}_{0}$ and $z^{\ddagger}_{0}$, respectively, 
			above the sources and the surface $\mathcal{S}''$, i.e., 
			with $z^{\dagger}_{0} < z^{t}_{0}$, $z^{\dagger}_{0} < z''_{0}$, 
			$z^{\ddagger}_{0} < z^{t}_{0}$ and $z^{\ddagger}_{0} < z''_{0}$. &
			$\mathbf{r}^{\dagger}$ and $\mathbf{r}^{\ddagger}$ &
			Generic points $\mathbf{r}^{\dagger} = (x^{\dagger}, y^{\dagger}, z^{\dagger}_{0})^{\top} \in \mathcal{S}^{\dagger}$ 
			and $\mathbf{r}^{\ddagger} = (x^{\ddagger}, y^{\ddagger}, z^{\ddagger}_{0})^{\top} \in \mathcal{S}^{\ddagger}$,
			both located above the sources and the surface $\mathcal{S}''$. \\
			$k_{x}$, $k_{y}$, $k$ and $\mathbf{k}$ &
			The first two variables represent the wavenumbers along the $x$ and $y$ directions, respectively, 
			while $k = \sqrt{k_{x}^{2} + k_{y}^{2}}$ and $\mathbf{k} = ( \mathtt{i} \, k_{x} \, , \, \mathtt{i} \, k_{y} \, , \, k)^{\top}$
			\citep{pedersen1978}, with $\mathtt{i}$ representing the imaginary unit. &
			$\mathcal{F} [ \cdot ]$ &
			Two-dimensional Fourier transform of a generic scalar function computed on a horizontal plane, with respect to
			the spatial coordinates $x$ and $y$. To simplify the notation, we conveniently omit the dependence of 
			$\mathcal{F} [ \cdot ]$ on the wavenumbers $k_{x}$ and $k_{y}$. \\
			$U$ and $U_{\varepsilon}$ & 
			Harmonic functions representing, respectively, the Newtonian potential produced by the true 
			sources (equation \ref{eq:potential-U-space}) and a smooth approximation for $U$ 
			(equation \ref{eq:potential-U-epsilon-space}). &
			$k_{\varepsilon}$ and $\mu_{\varepsilon}$ &
			Fixed truncating wavenumber and fixed positive regularization parameter defining the approximating potential 
			$U_{\varepsilon}$ according to equations \ref{eq:potential-U-epsilon-Fourier-dagger-1} and 
			\ref{eq:potential-U-epsilon-Fourier-dagger-2}, respectively. \\
			$d^{o}$ and $f$ &
			Harmonic functions representing potential fields produced by the true sources at points outside them 
			(equations \ref{eq:observed-data-space} and \ref{eq:transformed-field-space}). The first represents the observed potential field. &
			$p_{\varepsilon}$ &
			Continuous equivalent layer (Definition \ref{def:continuous-eqlayer}). \\
			$d_{\varepsilon}$ and $f_{\varepsilon}$ &
			Harmonic functions representing, respectively, the predicted and transformed potential fields produced by a 
			continuous equivalent layer $p_{\varepsilon}$ (equations \ref{eq:predicted-data-continuous-layer} and 
			\ref{eq:transformed-field-continuous-layer}). &
			$D_{d}$ and $\alpha_{d}$ &
			Space- and Fourier-domain operators that act on potentials $U$ and $U_{\varepsilon}$ to define the potential 
			fields $d^{o}$ and $d_{\varepsilon}$ (equations \ref{eq:observed-data-space}, \ref{eq:observed-data-Fourier} 
			and \ref{eq:predicted-data-continuous-layer-epsilon-Fourier}). \\
			$D_{a}$ and $\alpha_{a}$ &
			Space- and Fourier-domain operators that act on the inverse distance function $\Psi$ (equations \ref{eq:Psi} and \ref{eq:Psi-Fourier}) 
			to define the kernel $a$ (equations \ref{eq:predicted-data-continuous-layer} and \ref{eq:kernel-a-Fourier}). &
			$D_{f}$ and $\alpha_{f}$ &
			Space- and Fourier-domain operators that act on potentials $U$ and $U_{\varepsilon}$ to define the potential 
			fields $f$ and $f_{\varepsilon}$ (equations \ref{eq:transformed-field-space}, \ref{eq:transformed-field-Fourier} 
			and \ref{eq:transformed-field-continuous-layer-single-step}). \\
			$D_{t}$ and $\alpha_{t}$ &
			Space- and Fourier-domain operators that act on the inverse distance function $\Psi$ (equations \ref{eq:Psi} and \ref{eq:Psi-Fourier}) 
			to define the kernel $t$ (equations \ref{eq:transformed-field-continuous-layer}, \ref{eq:kernel-t-Fourier} and \ref{eq:alpha-t-generic}). &
			$\hat{\mathbf{u}}$, $\hat{\mathbf{v}}$, $\hat{\mathbf{h}}$, $\hat{\mathbf{x}}$, $\hat{\mathbf{y}}$ and $\hat{\mathbf{z}}$ &
			Constant unit vectors. The last three define the canonical basis (equation \ref{eq:canonical-base-vector}).
		\end{tabular}
	\end{center}
	\label{tbl:Table1}
}


\tabl{Table2}{Relevant potential fields for geophysical applications.
}{
	\begin{center}
		\renewcommand{\arraystretch}{0.8}
		\begin{tabular}{
				>{\small\centering}m{0.15\textwidth}
				>{\smaller\setlength{\baselineskip}{0.5\baselineskip}\centering}m{0.25\textwidth}
				>{\tiny\setlength{\baselineskip}{0.5\baselineskip}\raggedright\arraybackslash}m{0.60\textwidth}
			}
			{\small \textbf{Symbol}} & 
			{\small \textbf{Potential field}} & 
			{\small \textbf{Brief description}} \\
			\hline
			$U$ & 
			Newtonian potential &
			In gravity applications, $U$ (equation \ref{eq:potential-U-space}) represents the gravitational potential. 
			In magnetic applications, it represents an auxiliary potential whose first directional derivative along the 
			constant magnetization direction of the sources gives the magnetic scalar potential. \\
			$\tfrac{\partial}{\partial u} U(\mathbf{r})$ &
			Gravitational attraction component or pseudogravity anomaly &
			First derivative of $U$ along the Cartesian directions $u \in \{ x, y, z \}$ (equation \ref{eq:potential-U-1st-deriv-space}). 
			In gravity applications, the vertical component $\tfrac{\partial}{\partial z} U$ 
			is called \textit{gravity disturbance} \cite[e.g.,][ p. 93]{hofmann-wellenhof-moritz2005}. 
			The debate around the differences between gravity anomaly and gravity disturbance is beyond the scope of the 
			present work and has been carried about from a more geodetic than a geophysical point of view 
			\citep[e.g.,][]{li2001, fairhead2003, hackney-featherstone2003, hinze2005, vajda-etal2006, vajda-etal2007, vajda-etal2008}. 
			In magnetic applications, the vertical 
			component represents the \textit{pseudogravity anomaly} 
			\citep[][ p. 345]{blakely1996}. \\
			$\tfrac{\partial^{2}}{\partial u \, \partial v} U(\mathbf{r})$ &
			Gravitational tensor component or RTP anomaly &
			Second derivative of $U$ along the Cartesian directions $u, v \in \{ x, y, z \}$ 
			(equation \ref{eq:potential-U-2nd-deriv-space}). 
			In gravity applications, the second derivative represents a component of the 
			\textit{gravitational tensor} \citep[e.g.,][]{pedersen_rasmussen1990}.
			In magnetic applications, the particular second derivative 
			$\tfrac{\partial^{2}}{\partial z^{2}} U(\mathbf{r})$ along the vertical 
			direction represents the \textit{RTP anomaly}. 
			Differently from generic total-field anomalies 
			\cite[e.g.,][ p. 178]{blakely1996}, which have a ``skewed'' pattern 
			oscillating from negative to positive values and do not have their maxima placed 
			right above the magnetic sources, the RTP anomaly is easier to interpret 
			because it is predominantly positive right above the sources and
			decays to zero towards their boundaries. \citet{baranov1957} was the first
			to propose a method, called \textit{reduction to the pole} (RTP), to transform 
			generic total-field anomalies into this special anomaly. \\
			$\nabla U(\mathbf{r})^{\top} \hat{\mathbf{v}}$ &
			Magnetic scalar potential &
			Directional derivative of $U$ along the constant magnetization direction $\hat{\mathbf{v}}$ of the sources 
			(equation \ref{eq:potential-U-oblique-1st-deriv-space}), where $\nabla U(\mathbf{r})$ is the gradient of 
			$U$ (equation \ref{eq:potential-U-grad-space}). \\
			$\hat{\mathbf{u}}^{\top} \mathbf{H}_{U} \, \hat{\mathbf{v}}$ &
			Magnetic induction field component &
			Directional derivative of $U$ along $\hat{\mathbf{v}}$, the constant magnetization direction of the sources, 
			and $\hat{\mathbf{u}}$ (equation \ref{eq:potential-U-oblique-2nd-deriv-space}), where $\mathbf{H}_{U}$ is the 
			Hessian of $U$ (equation \ref{eq:potential-U-Hessian}). For the case in which 
			$\hat{\mathbf{u}}$ is $\hat{\mathbf{x}}$, $\hat{\mathbf{y}}$ or $\hat{\mathbf{z}}$ (equation \ref{eq:canonical-base-vector}),
			the directional derivative represents a Cartesian component of the magnetic induction field. 
			For the case in which $\hat{\mathbf{u}}$ is the constant direction of the main field, the directional derivative 
			represents the \textit{approximated total-field anomaly} \cite[e.g., ][ p. 179]{blakely1996}. \\
			$\| \nabla \left[\nabla U(\mathbf{r})^{\top} \hat{\mathbf{v}}\right] \|_{2}$ &
			Amplitude of the magnetic induction field &
			Square root of the sum of the squares of the Cartesian magnetic induction field components. This field, also known as 
			\textit{amplitude of the anomalous magnetic field} \citep{li_etal2010} or 
			\textit{amplitude of the magnetic anomaly vector} \citep{hidalgo_gato_etal2021, melo_etal2021}, has a weak dependence on the total-magnetization direction of the
			true sources. \\
		\end{tabular}
	\end{center}
	\label{tbl:Table2}
}
% Define figure directory
\renewcommand{\figdir}{figures}

\plot{Figure1}{width=0.8\textwidth}{
%\plot{}{width=0.8\textwidth}{
	{
		2D schematic representation of a vertical $yz$ plane crossing the sources (represented in dark gray),
		with $x$ axis pointing inward the page (represented by the ``$\otimes$" symbol).
		The sources are located in the subsurface (represented in light gray), over a 
		closed region $\overline{\mathcal{V}'}$ formed by $N_{\nu} = 2$ disjoint closed regions 
		$\overline{\mathcal{V}'_{\nu}}$, $\nu \in \{1, 2\}$, with top limited by the plane $z^{t}_{0}$. 
		The fine dashed lines represent two possible horizontal surfaces $\mathcal{S}''$, one located 
		at $z''_{0} < z^{t}_{0}$, above the sources, and another at $z''_{0} > z^{t}_{0}$, crossing the sources.
		The medium dashed line represents the surface containing the observation points 
		$\mathbf{r}$, all above $\mathcal{S}''$.
		The long dashed line represents a planar surface $\mathcal{S}^{\dagger}$ located 
		at $z^{\dagger}_{0} < z''_{0}$.
		Finally, the dashed-dot line represents a planar surface $\mathcal{S}^{\ddagger}$ located 
		at $z^{\ddagger}_{0} < z^{\dagger}_{0}$.
	}
	\label{fig:Figure1}
}



\clearpage

% Appendices
%\appendix
\append{Some useful upper bounds}
\label{app:upper-bounds}

We present here an upper bound for the squared absolute value of 
$\alpha(\ell,\hat{\mathbf{u}}) \, \alpha(m,\hat{\mathbf{v}}) \, \alpha(n,\hat{\mathbf{w}})$,
that is defined by equation \ref{eq:diff-operators-generic} using $\ell, m , n \in \{0, 1, 2\}$
and constant unit vectors $\hat{\mathbf{u}}$, $\hat{\mathbf{v}}$ and $\hat{\mathbf{w}}$. 
To do this, we first convert the wavenumber $k_{x}$ and $k_{y}$ to polar coordinates 
$k_{x} = k \, \cos\theta$, $k_{y} = k \, \sin\theta$ and $k = \sqrt{k_{x}^{2} + k_{y}^{2}}$
to show that
\begin{equation}
	\begin{split}
		\vert \alpha(\ell,\hat{\mathbf{u}}) \, \alpha(m,\hat{\mathbf{v}}) \, \alpha(n,\hat{\mathbf{w}}) \vert^{2} 
		&= \vert (\mathbf{k} \cdot \hat{\mathbf{u}})^{\ell} \, (\mathbf{k} \cdot \hat{\mathbf{v}})^{m} \, (\mathbf{k} \cdot \hat{\mathbf{w}})^{n} \vert^{2} \\
		&= \left[ (\mathbf{k} \cdot \hat{\mathbf{u}})^{2} \right]^{\ell} \left[ (\mathbf{k} \cdot \hat{\mathbf{v}})^{2} \right]^{m} 
		\left[ (\mathbf{k} \cdot \hat{\mathbf{w}})^{2} \right]^{n} \\
		&\le \left[ k_{x}^{2} + k_{y}^{2} + k^{2} \right]^{\ell} \left[ k_{x}^{2} + k_{y}^{2} + k^{2} \right]^{m} 
		\left[ k_{x}^{2} + k_{y}^{2} + k^{2} \right]^{n} \\
		&\le 3^{q} k^{2 \, q} \: , \quad q = \ell+m+n \: ,
		\quad \ell, m , n \in \{0, 1, 2\} \: .
	\end{split}
\label{eq:alpha-upper-bound}
\end{equation}

We are also interested in defining an upper bound for the absolute value
of $\mathcal{F} [ U(\mathbf{r}^{\dagger}) ]$ (equation \ref{eq:potential-U-Fourier}), 
i.e., for the amplitude spectrum of $U(\mathbf{r})$ (equation \ref{eq:potential-U-space}). 
To do this, we have adapted the approach presented by
\citet{bhattacharyya1967} to define an upper bound for the amplitude spectrum of 
$\tfrac{\partial}{\partial z} U(\mathbf{r})$. 
We first stress that 
$e^{- k \, (z' - z^{\dagger}_{0})} \le e^{- k \, (z^{t}_{0} - z^{\dagger}_{0})}$ because
$z^{t}_{0} \le z'$, where $z^{t}_{0}$ defines the top of the sources and $z'$ the vertical
coordinate of a point inside the sources. Using this identity into  
$\mathcal{F} [ U(\mathbf{r}^{\dagger}) ]$ (equation \ref{eq:potential-U-Fourier}), we obtain
\begin{equation}
	\vert \mathcal{F} [ U(\mathbf{r}^{\dagger}) ] \vert \le 2 \pi \kappa 
	\frac{e^{- k \, \Delta z}}{k}
	\int\limits_{z^{t}_{0}}^{\infty}
	\mathcal{F} [ \sigma(\mathbf{r}') ] \:
	\: dz' \: , \quad \Delta z = z^{t}_{0} - z^{\dagger}_{0} > 0 \: .
	\label{eq:potential-U-spectrum-bound1}
\end{equation}
Manipulating $\mathcal{F} [ \sigma(\mathbf{r}') ]$ (equation \ref{eq:sigma-Fourier}), 
we obtain
\begin{equation}
	\vert \mathcal{F} [ \sigma(\mathbf{r}') ] \vert \le
	\int\limits_{-\infty}^{\infty}\int\limits_{-\infty}^{\infty}
	\vert \sigma(\mathbf{r}') \vert 
	\: dx'dy'
	\label{eq:sigma-spectrum-bound}
\end{equation}
and, consequently, that
\begin{equation}
	2 \pi \kappa \, 
	\int\limits_{z^{t}_{0}}^{\infty}
	\vert \mathcal{F} [ \sigma(\mathbf{r}') ] \vert
	\: dz' \le \Omega \: ,
	\label{eq:Omega-bound}
\end{equation}
where
\begin{equation}
	\Omega = 2 \pi \kappa \, \iiint\limits_{\mathcal{V}}
	\vert \sigma(\mathbf{r}') \vert 
	\: d\mathcal{V}'
	\label{eq:Omega}
\end{equation}
is a constant depending only on the position and shape of the sources. 
Then, by combining equations 
\ref{eq:potential-U-spectrum-bound1}--\ref{eq:Omega},
we conclude that the amplitude spectrum of 
$U(\mathbf{r})$ (equation \ref{eq:potential-U-space}) on a horizontal plane 
$z^{\dagger}_{0} < z^{t}_{0}$ located above the sources satisfies 
\begin{equation}
	\vert \mathcal{F} [ U(\mathbf{r}^{\dagger}) ] \vert \le
	\frac{e^{- k \, \Delta z}}{k} \, \Omega \: ,
	\: \quad \Delta z = z^{t}_{0} - z^{\dagger}_{0} > 0 \: .
	\label{eq:potential-U-spectrum-bound}
\end{equation}
From the upper bound of $\vert \mathcal{F} [ U(\mathbf{r}^{\dagger}) ] \vert$
(equation \ref{eq:potential-U-spectrum-bound}) and $\mathcal{F} [ d^{o}(\mathbf{r}^{\dagger}) ]$
(equation \ref{eq:observed-data-Fourier}), we also deduce
\begin{equation}
	\vert \mathcal{F} [ d^{o}(\mathbf{r}^{\dagger}) ] \vert \le
	\vert \alpha_{d} \vert \, \Omega \, 
	\frac{e^{- k \, \Delta z}}{k} \: ,
	\: \quad \Delta z = z^{t}_{0} - z^{\dagger}_{0} > 0 \: .
	\label{eq:observed-data-spectrum-bound}
\end{equation}
Note that, for the case in which $\alpha_{d} = k$ (equation \ref{eq:observed-data-Fourier}), 
we obtain the same result presented by \citet{bhattacharyya1967} as the upper bound for the 
amplitude spectrum of $\tfrac{\partial}{\partial z} U(\mathbf{r})$, whereas
$\alpha_{d} = 1$ leads to the upper bound defined by equation \ref{eq:potential-U-spectrum-bound}.
Equations \ref{eq:potential-U-spectrum-bound} and \ref{eq:observed-data-spectrum-bound} 
show that the amplitude spectra 
$\vert \mathcal{F} [ U(\mathbf{r}^{\dagger}) ] \vert$ and 
$\vert \mathcal{F} [ d^{o}(\mathbf{r}^{\dagger}) ] \vert$ on a plane $z^{\dagger}_{0} < z^{t}_{0}$ 
above the sources 
(i) are finite at all wavenumbers $k = \sqrt{k_{x}^{2} + k_{y}^{2}}$ and (ii) tend
to zero as $k$ goes to infinite.
These results confirm the well-known fact that 
$\mathcal{F} [ U(\mathbf{r}^{\dagger}) ]$ and $\mathcal{F} [ d^{o}(\mathbf{r}^{\dagger}) ]$ 
(equations \ref{eq:potential-U-Fourier} and \ref{eq:observed-data-Fourier}) are square integrable.
\append{Phase spectrum of the potential $U(\mathbf{r})$}
\label{app:potential-U-min-phase}

Let the phase spectrum \cite[e.g.,][ p. 261]{blakely1996} of  
$\mathcal{F} [ U(\mathbf{r}^{\dagger}) ]$ (equation \ref{eq:potential-U-Fourier}) 
on a horizontal plane $z^{\dagger}_{0}$ above the sources be defined as follows:
\begin{equation}
	\Theta [ U(\mathbf{r}^{\dagger}) ] = \tan^{-1} \left( 
	\frac{\text{Im} \, \mathcal{F} [ U(\mathbf{r}^{\dagger}) ]}{\text{Re} \, \mathcal{F} [ U(\mathbf{r}^{\dagger}) ]} 
	\right) \: ,
	\label{eq:potential-U-phase-spectrum}
\end{equation}
where ``$\text{Re}$'' and ``$\text{Im}$'' denote the real and imaginary parts, respectively.


According to equation \ref{eq:potential-U-phase-spectrum}, the phase spectrum of $U(\mathbf{r}^{\dagger})$ (equation \ref{eq:potential-U-Fourier}) 
represents a weighted sum of the Fourier transforms $\mathcal{F} [ \sigma(\mathbf{r}') ]$
(equation \ref{eq:sigma-Fourier}) over all depth slices crossing the sources,
where the weights are real functions that exponentially decay with depth $z'$.
It means that $\Theta [ U(\mathbf{r}^{\dagger}) ]$ attributes more weight to the shallow than to the deeper parts of the sources. 
Because of that, the phase spectrum of $U(\mathbf{r}^{\dagger})$ approximates that of $\sigma(\mathbf{r}')$ and, consequently, 
$U(\mathbf{r}^{\dagger})$ has approximately the same horizontal spatial distribution of the (shallow parts of) true sources.
The same reasoning can be applied to the phase spectra of any vertical integral or vertical derivative of the potential $U$.

\clearpage

% References
\bibliographystyle{seg}  % style file is seg.bst
\bibliography{references}

\clearpage

\end{document}
